% =====================================================================
% Stage 3 -- Detailed Refinement.
%
% The decisions and verifiable requirements that follow from the
% high-level refinement: further architecture decisions (ADR) at a
% narrower scope than the foundational one, down to the one that fixes
% the eviction-reconstruction mechanism, and the concrete,
% testable storage requirements (FRE) derived from the
% standards-compliance goal into what the storage
% layer must guarantee across a store -> retrieve cycle.
% =====================================================================

% ---------------------------------------------------------------------
% Persistence & space efficiency: architecture decisions derived from
% the plain-file-system ADR, each satisfying the derived constraints
% (REQ-EFFICIENCY.*) and framework constraints (FRE-COMPAT.*)
% stated in the high-level refinement.
% ---------------------------------------------------------------------
\begin{requirement}{ADR-STORAGE.OBJECTFILE-01}[draft][designConstraint]{Combined metadata and body file}

  \localreqdiagram*{adr2}{\reqid}{\reqdesc}

  \textbf{Context:} Each cached object has a small block of metadata and a body
  of arbitrary length, unknown when writing begins
  (\reqref{FRE-COMPAT.UNSIZED-01}), persisted on a plain file system
  (\reqref{ADR-PERSISTENCE.MECHANISM-01}). Block-granular allocation means a
  separate file holding only the small metadata would waste most of a block
  per object (\reqref{REQ-EFFICIENCY.FRAGMENTATION-01}).

  \textbf{Decision:} Store each object's metadata and body together in a
  single growable file, in file order: the metadata written ahead of the
  body content, the body appended as its chunks arrive. A metadata update
  that later outgrows the space reserved ahead of the body withdraws the
  record, and the object is fetched again
  (\reqref{FDS-STORAGE.METAUPDATE-01}); moving that metadata after the body
  instead is deferred to a later version (\reqref{ADR-STORAGE.METAMOVE-01}).

  \textbf{Alternatives considered:}
  \begin{itemize}
    \item \emph{A separate file holding only the object's metadata}, body
      kept in its own file: keeps each region independently sized and
      simple to read in isolation, but the small, mostly-empty metadata
      file wastes most of a block per object regardless of the body's
      size (\reqref{REQ-EFFICIENCY.FRAGMENTATION-01}). Committing the two
      files as one atomic unit also loses the single-rename publication
      this design enables (\reqref{FDS-STORAGE.PUBLISH-01}): a plain
      rename is atomic per path, not across two, so publishing both
      together would need a per-object subdirectory renamed as a whole
      --- an extra inode and directory entry per object, working against
      the same efficiency goal --- or a separate commit marker. Nor does
      splitting the files cleanly avoid in-place metadata updates
      (\reqref{FDS-STORAGE.METAUPDATE-01}): replacing the metadata file
      wholesale on every revalidation avoids mutating it in place, but
      pays a fresh inode and directory entry each time, so a design that
      cares about that cost would still want to reserve headroom and
      update it in place --- the same capacity accounting and
      mid-write-race handling this design already needs, just for a
      smaller file.
    \item \emph{Metadata in the body file's extended attributes}: no second
      file, and a small attribute can live inside the inode itself, costing
      no extra block. Rejected: a header set has no size cap
      (\reqref{FRE-COMPLIANCE.HEADERS-01}), while attribute space is
      capped\footnote{\href{https://man7.org/linux/man-pages/man7/xattr.7.html}{Linux
      manual page: \texttt{xattr(7)}}: ``The VFS-imposed limits on attribute
      names and values are 255 bytes and 64 kB, respectively'', and on ext2,
      ext3, and ext4 ``the total bytes used by the names and values of all of
      a file's extended attributes must fit in a single filesystem block''.}
      --- so large header sets could not be stored at all, and typical
      header segments overflow the in-inode space into a block anyway;
      attribute writes are also separate from the file's data, adding a
      second update channel to keep consistent with the body.
    \item \emph{Metadata as a trailer after the body}: the metadata's size
      no longer has to be reserved ahead of a body of unknown length, and a
      growing update simply rewrites the tail. Rejected: every update
      changes the file's length while readers may be streaming the body's
      final bytes and validating the length they were promised, and every
      lookup must locate the trailer through a fixed footer before it can
      read anything; a header-first record confines updates to a
      fixed-position region instead.
    \item \emph{A sidecar index} --- one embedded key-value store holding
      every object's metadata, with bodies in their own files: small
      metadata records pack tightly, and updates become transactional.
      Rejected: publishing an object then spans two stores, the body file
      and the index entry, which can no longer be made visible by a single
      rename and need an ordering and recovery protocol of their own; and it
      reintroduces the embedded engine, with its cross-process locking,
      that \reqref{ADR-PERSISTENCE.MECHANISM-01} declined.
  \end{itemize}

  \textbf{Consequences:}
  \begin{itemize}
    \item[+] No block is wasted on a separate, mostly-empty metadata file.
    \item[+] One growable file naturally holds a body whose size is not
      known when writing begins, and a streaming partial reader can follow
      the same file the leader is still writing
      (\reqref{FRE-COMPAT.STREAMING-01}).
    \item[--] Metadata and body share a single file, so reading the metadata
      requires opening the object file rather than a smaller dedicated one.
    \item[--] The on-disk format must encode the metadata length or body offset
      so the two regions can be told apart on read.
    \item[--] Metadata that outgrows the region ahead of the body costs a
      fetch of the whole object from the origin
      (\reqref{ADR-STORAGE.METAMOVE-01}).
  \end{itemize}

  \reqdiaglink{overview}{\textit{View the domain overview.}}

  \reqexpect{refine}{FDS-STORAGE.PUBLISH-01}{Specify how object publication is
  made atomic under concurrent same-key writes and process crashes.}
  \reqexpect{refine}{FDS-STORAGE.METAUPDATE-01}{Specify how a metadata update
  that grows past its record's reserved capacity is handled.}
  \reqexpect{refine}{SDS-LAYOUT.OBJECTFILE-01}{Specify the concrete on-disk
  record format --- header, key, metadata region, and body placement ---
  that realises the combined metadata-and-body file.}

  \refine{ADR-PERSISTENCE.MECHANISM-01}
  \trace{FRE-COMPAT.STREAMING-01}
  \trace{REQ-EFFICIENCY.FRAGMENTATION-01}
\end{requirement}

\begin{requirement}{ADR-PERSISTENCE.DURABILITY-01}[draft][designConstraint]{Host-crash safety via sync and CRC}

  \localreqdiagram*{adrdurability01}{\reqid}{\reqdesc}

  \textbf{Context:} After the host crashes or loses power, the cache may lose
  objects but must never serve content that differs from what was stored,
  nor an object whose purge it acknowledged
  (\reqref{HLR-PERSISTENCE.CRASHSAFE-01}); the same bar applies to incomplete
  stores (\reqref{FRE-COMPLIANCE.INCOMPLETE-01}), acknowledged purges
  (\reqref{FRE-COMPLIANCE.PURGEDURABLE-01}), and applied metadata updates
  (\reqref{FRE-COMPLIANCE.METAPERSIST-01}). A plain file system
  (\reqref{ADR-PERSISTENCE.MECHANISM-01}) gives none of this for free. A
  \texttt{write} returns while the bytes are still only in the page cache;
  syncing a file does not persist the directory entry that names
  it;\footnote{\href{https://man7.org/linux/man-pages/man2/fsync.2.html}{Linux
  manual page: \texttt{fsync(2)}}: ``Calling \texttt{fsync()} does not
  necessarily ensure that the entry in the directory containing the file has
  also reached disk. For that an explicit \texttt{fsync()} on a file
  descriptor for the directory is also needed.''} and a rename can reach the
  journal before the renamed file's data does. ext4 guards only the
  \emph{replace}-via-rename pattern against that, forcing the new file's data
  out before such a rename commits;\footnote{\href{https://www.kernel.org/doc/html/latest/admin-guide/ext4.html}{Linux
  kernel documentation: ``ext4 General Information''}, \texttt{auto\_da\_alloc}:
  the replace-via-rename pattern is detected so that ``the data blocks of the
  new file are forced to disk before the \texttt{rename()} operation is
  committed'', avoiding the ``zero-length'' problem.} a first-time publication
  under a new name gets no such help. A synchronous flush costs roughly
  1--5~ms on consumer SSDs and under 0.1~ms on drives with power-loss
  protection.\footnote{\href{https://smalldatum.blogspot.com/2026/01/ssds-power-loss-protection-and-fsync.html}{Mark
  Callaghan, ``Small Datum'': ``SSDs, power loss protection and fsync
  latency''}.}

  \textbf{Decision:} Synchronise exactly the writes whose loss could make the
  cache serve wrong or purged content, and detect the rest.
  \begin{itemize}
    \item \emph{Publication:} \texttt{fdatasync} a store's temporary ---
      after its last body byte and its patched header --- before renaming it
      to the canonical name. The directory is not synced after the rename,
      so a host crash may lose a just-published object (the key then keeps
      its previous record, or misses), but the name can never refer to
      data that is not on stable storage. This is the \emph{strict}
      configuration, the default: \reqref{ADR-PERSISTENCE.UNSYNCED-01}
      adds one that skips the synchronisation and discards the cache
      after a host restart instead.
    \item \emph{Invalidating purge:} \texttt{fsync} the shard directory
      after the \texttt{unlink}, or after finding nothing to unlink, before
      reporting the outcome. A purge made to evict is not synchronised: a
      host crash that undoes it brings back a valid object the cache chose
      to drop.
    \item \emph{Metadata update, in place or a withdrawal:} not synced.
      After a host crash the region holds the old metadata, the new, or a
      torn mixture, and a withdrawn record may be back as it was; the
      metadata checksum rejects the mixture, so the object is either served
      with one complete version or reported as a miss.
    \item \emph{Body checksum:} each record carries a CRC-32C of its body,
      computed while the body streams in and recorded at completion. Serving
      a hit does not check it. A build can leave it out (the Cargo feature
      \texttt{body-crc}, on by default; without it a record carries zero),
      and a build with the \texttt{body-crc-verify} feature can switch on, in
      its configuration, a read-time check that logs a mismatch once a read
      has delivered the whole body, and still ends the read normally.
  \end{itemize}
  With synchronised publication, on a storage stack that honours flush
  requests as \reqref{HLR-PERSISTENCE.CRASHSAFE-01} assumes, a host crash
  cannot leave a published name over wrong bytes. Integrity at rest --- media
  corruption, or a device that acknowledges a flush it has not made --- is
  the storage layer's: a data-checksumming file system (ZFS, btrfs,
  bcachefs) or \texttt{dm-integrity} turns a corrupted block into a read
  error, which ends the read with an error in place of the corrupted bytes.
  The header and metadata checksums stay: they detect a torn in-place
  metadata update, and a lookup racing one, which no storage layer can.

  \textbf{Alternatives considered:}
  \begin{itemize}
    \item \emph{Process-crash safety only} (no synchronisation): cheapest,
      and sufficient for \reqref{HLR-PERSISTENCE.DURABLE-01}, since a process
      crash leaves the page cache intact. Rejected: after a host crash a
      first-time rename can commit ahead of its data, leaving a canonical name
      over a zero-filled or short body whose header still passes validation,
      and an unsynced purge \texttt{unlink} can be lost, bringing a purged
      object back --- both forbidden by \reqref{HLR-PERSISTENCE.CRASHSAFE-01}.
      \reqref{ADR-PERSISTENCE.UNSYNCED-01} offers it as an option, made safe
      by discarding the cache after a host restart.
    \item \emph{Checksum only, without synchronisation}: detects a corrupted
      body, but only once a reader reaches its end, after the corrupted bytes
      have already gone to the client; and it does nothing for a lost purge.
      Rejected.
    \item \emph{Check the body checksum on every whole-body read}, ending a
      read whose body does not match with an error instead of a clean end of
      body. Rejected: the check can only run after the last byte, and
      \texttt{pingora-cache} sends each chunk downstream as the backend
      returns it, so a client given a \texttt{Content-Length} has the whole
      corrupted body by then and the error can only cut the connection. On
      the storage stack \reqref{HLR-PERSISTENCE.CRASHSAFE-01} assumes it never
      fires, yet it costs a CRC-32C pass per read, about as much CPU as the
      read itself when the body is in the page cache (13.3 against 12.3
      microseconds per 64~KiB chunk, measured on an 8-core host).
    \item \emph{Full host-crash durability}: additionally sync the directory
      after every rename and sync every in-place metadata update, so no
      completed store or update is ever lost. Rejected: it roughly doubles the
      synchronous cost of a store and adds one to every revalidation, to
      prevent only losses \reqref{HLR-PERSISTENCE.CRASHSAFE-01} explicitly
      accepts --- the kind of cost \reqref{HLR-SCOPE.V1-01} asks not to pay
      without a requirement behind it.
  \end{itemize}

  \textbf{Consequences:}
  \begin{itemize}
    \item[+] A host crash can lose recently stored objects or recent metadata
      updates, but can neither corrupt what is served nor resurrect an object
      an invalidation removed; an eviction it undoes brings back a valid
      object.
    \item[+] Process crashes and orderly restarts lose nothing that had been
      published (\reqref{HLR-PERSISTENCE.DURABLE-01}).
    \item[--] One \texttt{fdatasync} per completed store in the strict
      configuration, paid when the store finishes, and one directory
      \texttt{fsync} per invalidating purge; evictions pay none.
    \item[+] Serving a hit computes no checksum over the body.
    \item[--] The record header carries a body-checksum field, and every
      store computes a CRC-32C over its body, unless the build leaves it out.
    \item[--] Serving a hit does not detect a body corrupted at rest. That
      takes a storage layer that verifies data as it reads it, or the
      optional read-time check, which only logs the mismatch after the bytes
      have been sent.
    \item[--] The guarantees assume the storage stack honours flush requests.
      A device or virtualisation layer that acknowledges a flush it has not
      made can lose a purge, or leave a published name over wrong bytes,
      undetected.
  \end{itemize}

  \reqdiaglink{overview}{\textit{View the domain overview.}}

  \reqexpect{refine}{FDS-STORAGE.PUBLISH-01}{Specify the fdatasync of the
  temporary before its rename, and how a failed sync is surfaced.}
  \reqexpect{refine}{FDS-STORAGE.PURGE-01}{Specify the fsync of the shard
  directory before an invalidating purge is reported, and that a purge made
  to evict is not synchronised.}
  \reqexpect{refine}{FDS-STORAGE.METAUPDATE-01}{Specify that a metadata
  update, written in place or withdrawing the record, is not synced.}
  \reqexpect{refine}{SDS-LAYOUT.OBJECTFILE-01}{Add the body CRC-32C field
  and how it is computed and recorded, and specify the optional read-time
  check that logs a mismatch at the end of a whole-body read.}

  \refine{ADR-PERSISTENCE.MECHANISM-01}
  \satisfy{HLR-PERSISTENCE.CRASHSAFE-01}
  \trace{FRE-COMPLIANCE.INCOMPLETE-01}
  \trace{FRE-COMPLIANCE.PURGEDURABLE-01}
  \trace{FRE-COMPLIANCE.METAPERSIST-01}
\end{requirement}

\begin{requirement}{ADR-PERSISTENCE.UNSYNCED-01}[draft][designConstraint]{Unsynced stores, optional discard}

  \localreqdiagram*{adrunsynced01}{\reqid}{\reqdesc}

  \textbf{Context:} \reqref{ADR-PERSISTENCE.DURABILITY-01} synchronises every
  store's temporary before renaming it into place, so that a host crash can
  never leave a published name over data that is not on stable storage.
  Measured on 22~October~2026 with the benchmark of
  \reqref{TC-PERFORMANCE.STORERATE-01}, on a consumer SSD under ext4, each
  \texttt{fdatasync} took 16--25~ms with 4 to 64 stores syncing at once, and
  stores completed only as fast as their syncs: 2{,}591 and 1{,}757 stores per
  second for 4~KiB and 64~KiB bodies, 15--16\,\% of \texttt{MemCache}.
  Without the sync, the same stores ran at 7{,}966 and 5{,}427 per second,
  46--49\,\%. \reqref{HLR-PERFORMANCE.STORERATE-01} asks for 40\,\% where the
  integrator accepts losing stored objects in a host crash;
  \reqref{HLR-PERSISTENCE.DURABLE-01} still holds there, and
  \reqref{HLR-PERSISTENCE.CRASHSAFE-01} as far as this decision says.

  Without the sync, a host crash can leave a published name over a body that
  never reached the device. The read validation of
  \reqref{SDS-LAYOUT.OBJECTFILE-01} catches the common case: it compares a
  record's file length with the record length its header gives, and a file
  system that records a file's length only once its data is written --- ext4
  in its default \texttt{data=ordered} mode, XFS, btrfs --- leaves such a
  record short or empty, which is then a miss. It cannot catch a body of the
  right length that reads back as zeros, which a file system that records the
  length first can leave (ext4 with \texttt{data=writeback}), and which is
  why \reqref{ADR-PERSISTENCE.DURABILITY-01} rejects process-crash safety
  alone. What the cache can tell, on any file system, is whether the host has
  restarted since it last published a store without synchronising it: Linux
  gives every boot a random identity, its boot
  ID.\footnote{\href{https://man7.org/linux/man-pages/man4/random.4.html}{Linux
  manual page: \texttt{random(4)}}, \texttt{/proc/sys/kernel/random/boot\_id}.}

  \textbf{Decision:} A configuration option chooses how a store becomes
  durable. \emph{Strict}, the default, is \reqref{ADR-PERSISTENCE.DURABILITY-01}
  as it stands. \emph{Unsynchronised} renames each store into place without
  the \texttt{fdatasync}. After a host crash it keeps every record the file
  system recovered: the read validation turns each one left short or empty
  into a miss, and the restart scan reclaims it
  (\reqref{FDS-STORAGE.RECLAIM-01}). On a file system that records a file's
  length only once its data is written, the crash costs only the stores whose
  data had not reached the device; on one that records the length first, a
  body of zeros could be served, a risk the integrator who chooses this
  configuration accepts. A second switch, off by default and on Linux only,
  removes that risk on any file system: the \emph{discard} makes the cache
  start empty after any host restart that may have lost what it published
  unsynchronised.
  \begin{itemize}
    \item Before it publishes a store unsynchronised with the discard on, a
      backend instance that has no marker writes one, durably, beside its
      writer lease (\reqref{ADR-STORAGE.WRITERID-01}): the boot ID it runs
      under.
    \item When it becomes read-only (\reqref{FDS-COMPAT.READONLY-01}), or the
      integrator flushes it, the backend flushes the file system that holds
      the cache and then removes, durably, its own marker and those of
      instances no longer running: everything they published is then on
      stable storage. A later unsynchronised store, with the discard on,
      writes the marker again.
    \item At construction, whatever its configuration, a marker naming a boot
      other than the current one means the host restarted while stores
      published unsynchronised could still be lost. The backend then discards
      every record and temporary before it serves anything, and only then
      removes the markers. Markers naming the current boot mean only a
      process restarted, and the page cache kept everything: nothing is
      discarded.
    \item Invalidating purges synchronise their shard directory in both
      configurations: a purge acknowledged while no marker exists must
      survive a host crash that the next start will not detect.
  \end{itemize}

  Under \reqref{HLR-SCOPE.V1-01} this takes the exception: the mechanisms it
  adds --- the option, the markers, the flush and the discard --- raised the
  store rates \reqref{TC-PERFORMANCE.STORERATE-01} measures 3.1 times, from
  15--16\,\% to 46--49\,\% of \texttt{MemCache}.

  \textbf{Alternatives considered:}
  \begin{itemize}
    \item \emph{The discard on by default.} Not adopted: on the file systems
      most deployments use --- ext4 in its default mode, XFS, btrfs --- the
      read validation already turns what a crash leaves behind into misses,
      so a default discard would empty the cache after every host crash to
      guard against file systems that record a file's length before its
      data.
    \item \emph{Check every record's body checksum after a host restart},
      keeping the records that pass. Rejected: the cache could serve nothing
      before reading every byte it holds, a wait that grows with the cache
      (\reqref{HLR-PERSISTENCE.CAPACITY-01}); and a build that leaves the
      checksum out has none to check.
    \item \emph{Periodic sync points}: flush the file system every few
      seconds, number each record with the last completed sync point, and
      after a host crash discard only the records numbered after it. It loses
      seconds of stores rather than the cache. Not adopted: it adds a field to
      the record format, a background task, and a check of every record at
      restart --- a candidate for a later version, once a measured figure
      shows that the empty restart matters.
    \item \emph{Asynchronous publication}: report a store complete at once,
      then synchronise and rename it in the background, keeping every
      guarantee of \reqref{ADR-PERSISTENCE.DURABILITY-01}. Rejected: a store
      reported complete is then not yet visible by name, which changes what a
      purge, the eviction manager's admission and the read-only drain see
      (\reqref{FDS-STORAGE.PUBLISH-01}), each to be specified again.
    \item \emph{Group commit}: one flush for a batch of stores. Rejected: each
      store still waits for a flush, and ext4 already commits concurrent
      \texttt{fdatasync} calls together, which the measurements above
      include.
  \end{itemize}

  \textbf{Consequences:}
  \begin{itemize}
    \item[+] A store no longer waits for the device: 3.1 times the store rate
      on a consumer SSD.
    \item[+] \reqref{HLR-PERSISTENCE.DURABLE-01} holds in both
      configurations: a process exit loses nothing.
    \item[+] On a file system that records a file's length only once its
      data is written, a host crash costs only the stores whose data had
      not reached the device.
    \item[--] Unsynchronised, \reqref{HLR-PERSISTENCE.CRASHSAFE-01} holds
      only on such a file system unless the discard is on: elsewhere, a host
      crash can leave a body of zeros to be served.
    \item[+] With the discard on, \reqref{HLR-PERSISTENCE.CRASHSAFE-01}
      holds on any file system.
    \item[--] With the discard on, a host crash, or any host restart the
      integrator did not prepare by flushing the backend, empties the cache,
      which refills from the origin.
    \item[--] The discard is Linux only, for its boot ID and
      \texttt{syncfs}; elsewhere it cannot be switched on.
    \item[--] Every backend reads the markers at construction, whatever its
      configuration, so that a cache last written with the discard on is
      never trusted after a host crash.
  \end{itemize}

  \reqdiaglink{overview}{\textit{View the domain overview.}}

  \reqexpect{refine}{FDS-PERSISTENCE.CRASHDISCARD-01}{Specify the options, the
  boot identity, when a marker is written and removed, the check at
  construction, how the discard keeps every record from being served, and
  what protects the records when the discard is off.}

  \refine{ADR-PERSISTENCE.DURABILITY-01}
  \satisfy{HLR-PERFORMANCE.STORERATE-01}
  \satisfy{HLR-PERSISTENCE.CRASHSAFE-01}
  \trace{HLR-PERSISTENCE.DURABLE-01}
  \trace{HLR-SCOPE.V1-01}
\end{requirement}

\begin{requirement}{FDS-PERSISTENCE.CRASHDISCARD-01}[draft][functionalRequirement]{Discard after a host restart}

  \localreqdiagram*{fdscrashdiscard01}{\reqid}{\reqdesc}

  IF the backend, at construction, finds an unsynchronised-store marker naming
  a boot other than the host's current one, THEN the backend shall make every
  record and temporary in the cache unreachable, by lookups and by the restart
  scan, before it answers its first lookup
  (\reqref{ADR-PERSISTENCE.UNSYNCED-01}).

  \textbf{The options.} \texttt{FsCacheConfig} carries the durability of
  stores: \texttt{Strict}, the default, or \texttt{NoStoreSync}. With
  \texttt{NoStoreSync}, \texttt{finish} renames a store into place without
  synchronising its temporary first (\reqref{FDS-STORAGE.PUBLISH-01});
  everything else is as with \texttt{Strict}, the invalidating purge's
  directory sync included (\reqref{FDS-STORAGE.PURGE-01}). A second switch,
  off by default, turns on the discard after a host restart. Off, the
  backend writes no marker, and after a host crash it keeps every record the
  file system recovered: the read validation of
  \reqref{SDS-LAYOUT.OBJECTFILE-01}, which compares a published record's
  file length with the record length its header gives, reports each record
  left short or empty as a miss, and the restart scan reclaims it
  (\reqref{FDS-STORAGE.RECLAIM-01}). The discard is available on Linux
  only: elsewhere, constructing a backend with it switched on fails with an
  error that says so.

  \textbf{The boot identity.} The current boot's identity is the content of
  \texttt{/proc/sys/kernel/random/boot\_id}. IF the backend cannot read it,
  THEN constructing a backend with the discard switched on fails with an
  error naming the file, and any other backend that finds markers discards
  the cache, since it cannot tell their boot from the current one.

  \textbf{The marker.} Each backend instance with the discard on has its own
  marker beside its writer lease (\reqref{FDS-STORAGE.PUBLISH-01}),
  \texttt{writers/<pid>-<nonce>.unsynced}, holding the boot identity. Before
  a \texttt{NoStoreSync} store's rename, if the instance's marker does not
  exist, the backend writes it --- a temporary, synchronised, renamed into
  place, and the \texttt{writers} directory synchronised --- once, however
  many stores reach that point together; the rename waits for it.

  \textbf{Flushing.} The backend shall expose a \texttt{flush} operation,
  which \texttt{set\_read\_only} (\reqref{FDS-COMPAT.READONLY-01}) also
  performs. It flushes the file system that holds the cache with
  \texttt{syncfs}, then removes the instance's own marker and the markers of
  instances whose lease is no longer held --- the same flush made their data
  durable --- and synchronises the \texttt{writers} directory. A
  \texttt{NoStoreSync} store published afterwards, with the discard on,
  writes the marker again first. With the discard on, an integrator that
  shuts its server down in order calls \texttt{flush} once the server has
  stopped; one that does not loses the cache at the next host restart.

  \textbf{At construction.} Before the shard-directory cache is built
  (\reqref{FDS-LAYOUT.STARTUP-01}), the backend reads every marker, whatever
  its configuration. If any
  names a boot other than the current one, it moves every first-level shard
  directory into a new directory, \texttt{discarded-<nonce>}, under the cache
  root, synchronises the root, and only then removes every marker and
  synchronises the \texttt{writers} directory: a crash during the discard
  leaves the markers, and with them the discard to be done again. The
  shard-directory cache then starts from new, empty first-level directories,
  and every \texttt{discarded-<nonce>} directory --- this start's, or one a
  crash left behind --- is removed on the blocking pool in the background,
  outside the tree the restart scan (\reqref{FDS-EVICTION.STARTUP-01}) walks.
  Markers that all name the current boot are left in place: only a process
  restarted, and the page cache kept everything.

  \reqdiaglink{overview}{\textit{View the domain overview.}}

  \reqcover{verify}{TC}

  \refine{ADR-PERSISTENCE.UNSYNCED-01}
  \satisfy{HLR-PERSISTENCE.CRASHSAFE-01}
  \trace{FDS-STORAGE.PUBLISH-01}
  \trace{FDS-COMPAT.READONLY-01}
  \trace{FDS-LAYOUT.STARTUP-01}
  \trace{FDS-EVICTION.STARTUP-01}
  \trace{FDS-STORAGE.RECLAIM-01}
\end{requirement}

\begin{requirement}{ADR-STORAGE.ATOMICPUBLISH-01}[draft][designConstraint]{Publish by same-directory rename}

  \localreqdiagram*{adratomicpublish01}{\reqid}{\reqdesc}

  \textbf{Context:} A store's record grows over time
  (\reqref{FRE-COMPAT.UNSIZED-01}) in a single file
  (\reqref{ADR-STORAGE.OBJECTFILE-01}), yet every write must be durably
  published as a complete object or not at all, including when several
  same-key writers race (\reqref{FRE-COMPAT.CONCURRENT-01}), a store is
  abandoned half-way (\reqref{FRE-COMPLIANCE.INCOMPLETE-01}), or a new
  response replaces an old one that readers still hold
  (\reqref{FRE-COMPLIANCE.REPLACE-01},
  \reqref{FRE-COMPLIANCE.READERISOLATION-01}). The question is how a completed
  record becomes visible by name --- to a scan, a restart, or a peer process
  --- in one indivisible step.

  \textbf{Decision:} Write each store to its own temporary file in the
  \emph{same directory} as the object's canonical name, and publish it with a
  single POSIX \texttt{rename} over that name once the store is complete.
  POSIX makes \texttt{rename} atomic between two paths on the same file
  system, replacing any existing target in the same step, while a descriptor
  already open on the old file keeps reading it. Keeping the temporary in the
  canonical name's directory makes that precondition structural rather than a
  deployment convention. The guarantee assumes POSIX \texttt{rename}
  semantics, matching the platform focus of \texttt{pingora} itself: Linux is its
  tier-1 environment, and Windows support is preliminary, community
  best-effort only.\footnote{The \texttt{pingora} README states this
  directly: \pgsrc{README.md}{53}{``Linux is our tier 1 environment and main
  focus.''}, and \pgsrc{README.md}{57}{``Windows support is preliminary by
  community's best effort only.''}} Windows is out of scope for this
  guarantee.

  \textbf{Alternatives considered:}
  \begin{itemize}
    \item \emph{Write in place under the canonical name, with a ``complete''
      flag in the header}: no temporary name and nothing to rename. Rejected:
      a peer process or a restart sees the partial record under the real
      name and must trust the flag; a replacement overwrites the old object
      in place while readers may still be streaming it; and a crash mid-write
      leaves a torn record where a complete one used to be.
    \item \emph{Stage the temporary in a separate scratch directory}:
      simpler cleanup, since every leftover sits in one place. Rejected:
      \texttt{rename} is atomic only within one file system, and one attempted
      across devices fails outright (\texttt{EXDEV}); a copy-then-delete
      fallback would reopen exactly the incomplete-publication window this
      decision exists to close.
    \item \emph{An unnamed temporary (\texttt{O\_TMPFILE}) published with
      \texttt{linkat}}: the kernel discards an unlinked temporary when its last
      descriptor closes, so a crash leaves no leftovers at
      all.\footnote{\href{https://man7.org/linux/man-pages/man2/open.2.html}{Linux
      manual page: \texttt{open(2)}}, \texttt{O\_TMPFILE}: ``Anything written
      to the resulting file will be lost when the last file descriptor is
      closed, unless the file is given a name.''} Rejected for this version:
      it is Linux-specific and depends on file-system support, and a link
      never replaces an existing
      name,\footnote{\href{https://man7.org/linux/man-pages/man2/link.2.html}{Linux
      manual page: \texttt{link(2)}}: ``If \emph{newpath} exists, it will not
      be overwritten.''} so a replacement would still need a named
      intermediate and a rename. The leftovers it would avoid are already
      reclaimed by a scan the backend runs anyway
      (\reqref{ADR-STORAGE.RECLAIM-01}).
    \item \emph{Windows \texttt{MoveFileEx}}: by default it refuses to
      overwrite an existing destination, and even with
      \texttt{MOVEFILE\_REPLACE\_EXISTING} the replace can fail with a sharing
      violation while another handle holds the destination open --- a failure
      mode POSIX \texttt{rename} does not have. Out of scope with Windows.
  \end{itemize}

  \textbf{Consequences:}
  \begin{itemize}
    \item[+] Publication, replacement, and the concurrent-writer race all
      reduce to one atomic step: the last rename wins, and readers of the old
      object are undisturbed.
    \item[+] Atomicity follows from placement alone; no deployment setting
      can break it.
    \item[--] A crash leaves named temporaries behind, which a later scan must
      reclaim (\reqref{ADR-STORAGE.RECLAIM-01}), and those temporaries share a
      directory with canonical names (\reqref{SDS-LAYOUT.SHARDING-01}).
    \item[--] POSIX-only.
  \end{itemize}

  \reqdiaglink{overview}{\textit{View the domain overview.}}

  \reqexpect{refine}{FDS-STORAGE.PUBLISH-01}{Specify the temporary's naming
  and placement, and the completion sequence that ends in the rename.}

  \refine{ADR-STORAGE.OBJECTFILE-01}
  \satisfy{FRE-COMPAT.CONCURRENT-01}
  \satisfy{FRE-COMPLIANCE.INCOMPLETE-01}
  \satisfy{FRE-COMPLIANCE.REPLACE-01}
  \satisfy{FRE-COMPLIANCE.READERISOLATION-01}
  \trace{ADR-STORAGE.RECLAIM-01}
\end{requirement}

\begin{requirement}{ADR-STORAGE.WRITERID-01}[draft][designConstraint]{Writer identity in temp names}

  \localreqdiagram*{adrwriterid01}{\reqid}{\reqdesc}

  \textbf{Context:} Each store writes a named temporary
  (\reqref{ADR-STORAGE.ATOMICPUBLISH-01}) that must never collide with
  another writer's, even across a restart, and a crash's leftover
  temporaries must be reclaimed by the first process start after their writer
  stopped (\reqref{FRE-COMPAT.STARTRECLAIM-01}) --- without ever removing
  one a live writer, such as a hot-upgrade peer
  (\reqref{FRE-COMPAT.HOTUPGRADE-01}), still depends on. A reclaiming scan
  therefore has to tell, from a temporary's name alone, whether its writer
  could still be running. A process ID alone cannot: IDs are reused, and a
  cache process running as PID~1 in a container gets the same ID on every
  start.

  \textbf{Decision:} Give each backend instance a \emph{lease}: a file under
  the cache root, \texttt{writers/<pid>-<nonce>.lock}, locked with
  \texttt{flock} for as long as the instance lives, where \texttt{<nonce>} is
  a random 64-bit value. Name each temporary after it: the object's hash, the
  lease's identity \texttt{<pid>-<nonce>}, the instance's own store counter,
  then the transient extension. A temporary's writer is considered possibly running
  exactly when its lease exists and its lock is held; the kernel releases a
  lock when its holder's descriptors close, a crash
  included.\footnote{\href{https://man7.org/linux/man-pages/man2/flock.2.html}{Linux
  manual page: \texttt{flock(2)}}: a lock created by \texttt{flock} is
  associated with an open file description, and is released by an explicit
  unlock ``or when all such file descriptors have been closed''.} The nonce
  makes the identity unique; the PID only tells an operator which process
  holds the lease. A lease is created under a staging name, locked, and only
  then renamed into place, so a scan never finds an unlocked lease under its
  final name while its writer lives; the restart-time scan removes the
  leases of writers that have stopped. A temporary is created exclusively,
  so a name collision is detected rather than silently reusing another
  writer's file. All processes sharing one cache directory must run on one
  host, with the directory on a local file system, so that each sees the
  others' locks; this is a deployment constraint the deliverable documents.

  \textbf{Alternatives considered:}
  \begin{itemize}
    \item \emph{PID only, tested with \texttt{kill(pid, 0)}}: the earliest
      design. Rejected: a reused PID makes a stopped writer look alive, and
      with PID~1 in containers every crash leftover looks alive forever and
      collides with the next process's temporaries once its counter restarts;
      \texttt{EPERM} also leaves liveness ambiguous.
    \item \emph{The writer's process instance}: a token of the host's boot
      identifier, the PID, and the process's start time, which a scan
      compares with the process's \texttt{/proc} entry. The previous
      decision. Rejected: it depends on Linux's \texttt{/proc} and boot
      identifier, and every process sharing a directory must run in one PID
      namespace, which excludes containers on one host sharing a volume.
    \item \emph{A lock on each temporary}: the writer locks its temporary for
      the life of the store and the scan tests that lock --- no new on-disk
      structure. Rejected: a temporary exists unlocked between its creation
      and its lock, which only a minimum-age rule could cover, and every
      store pays for one more lock; a lease is locked once, before its
      instance creates any temporary.
    \item \emph{A random name, reclaimed once older than a threshold}:
      rejected --- it ties correctness to a timeout that a long, slow upstream
      body can always exceed.
  \end{itemize}

  \textbf{Consequences:}
  \begin{itemize}
    \item[+] A stopped writer's leftovers --- a crash's, a killed process's, or
      an earlier boot's --- are recognised at the first scan, and no reused
      PID can make a stopped writer look alive.
    \item[+] Two writers can never share a temporary name.
    \item[+] No dependency on \texttt{/proc}, a boot identifier, or a PID
      namespace: processes in separate PID namespaces on one host may share a
      cache directory.
    \item[--] A new on-disk structure (\reqref{HLR-SCOPE.V1-01}): one empty
      lease per running instance, in a directory of its own, which the
      restart-time scan cleans up.
    \item[--] A lease outlives its instance until the next restart-time scan.
    \item[--] The cache directory must be on a local file system. On NFS,
      Linux emulates \texttt{flock} with POSIX record locks, which never
      conflict within one process and are all released when the process
      closes any descriptor of the file: a scan would take its own
      process's lease for a stopped writer's.
    \item[--] Temporary names grow longer.
  \end{itemize}

  \reqdiaglink{overview}{\textit{View the domain overview.}}

  \reqexpect{refine}{FDS-STORAGE.PUBLISH-01}{Specify the temporary name
  format and exclusive creation, including what happens on a collision.}
  \reqexpect{refine}{FDS-STORAGE.RECLAIM-01}{Specify the liveness test the
  scan applies to a temporary's writer (its lease's lock), the removal of
  stopped writers' leases, and the deployment constraint.}

  \refine{ADR-STORAGE.ATOMICPUBLISH-01}
  \satisfy{FRE-COMPAT.STARTRECLAIM-01}
  \trace{FRE-COMPAT.CONCURRENT-01}
  \trace{FRE-COMPAT.HOTUPGRADE-01}
  \trace{ADR-STORAGE.RECLAIM-01}
\end{requirement}

\begin{requirement}{FDS-STORAGE.PUBLISH-01}[draft][functionalRequirement]{Atomic publication via rename}

  \localreqdiagram*{fdspublish01}{\reqid}{\reqdesc}

  \textbf{The writer's lease.} When the backend is constructed, it shall
  create its lease (\reqref{ADR-STORAGE.WRITERID-01}) under the cache root
  as \texttt{writers/<pid>-<nonce>.lock.new}, exclusively, take an exclusive
  \texttt{flock} on it, and rename it to \texttt{writers/<pid>-<nonce>.lock},
  choosing a new nonce if the rename finds the file gone. It holds the lock,
  with the file open, for as long as it lives.

  \textbf{Opening a store.} When a store begins, the backend shall create the
  store's \emph{per-writer temporary} exclusively --- failing rather than
  opening an existing file --- in the same directory as the object's
  canonical name (\reqref{SDS-LAYOUT.SHARDING-01}), so the final rename is
  atomic by construction (\reqref{ADR-STORAGE.ATOMICPUBLISH-01}). The
  canonical name is the object's name as \reqref{SDS-LAYOUT.SHARDING-01}
  defines it: the hex-encoded combined hash of the cache
  key,\footnote{\pgsrc{pingora-cache/src/key.rs}{90}{\texttt{CacheHashKey::combined}}.}
  or, for a key with a user tag, a hash of that hash and the tag. The
  temporary's name appends the writer identity
  \reqref{ADR-STORAGE.WRITERID-01} decides, the instance's own write
  counter, and the transient extension:
  \texttt{<hash>.<pid>-<nonce>.<counter>.tmp}. Here
  \texttt{<pid>-<nonce>} names the lease of the backend instance writing it
  --- its process ID in decimal and a random 64-bit value in sixteen
  hexadecimal digits --- and \texttt{<counter>}, in decimal, is incremented
  for every store the instance opens. Each instance's identity keeps its
  names apart from every other instance's, and the counter keeps the
  instance's own writers of one key apart --- the cache lock is optional, and
  its age timeout lets a new writer start while a previous one may still be
  writing (\reqref{FRE-COMPAT.CONCURRENT-01}) --- just as the stock
  \texttt{MemCache} keys its in-flight writes by a per-writer
  identifier.\footnote{\pgsrc{pingora-cache/src/memory.rs}{364}{\texttt{MemCache::get\_miss\_handler}}
  assigns each new write its own incrementing id and inserts it into the
  per-key map alongside any other writes already in flight.} If exclusive
  creation reports that the name exists, the backend shall advance the
  counter and retry.

  The backend shall then write the record header, key, and metadata region
  (\reqref{SDS-LAYOUT.OBJECTFILE-01}), and append each body chunk as
  \texttt{write\_body} delivers it, folding it into a running CRC-32C unless
  the build leaves out the \texttt{body-crc} feature
  (\reqref{ADR-PERSISTENCE.DURABILITY-01}). While the temporary grows it is
  also the read source for streaming partial hits, as
  \reqref{FDS-STORAGE.PARTIAL-01} specifies.

  \textbf{Write-through.} The backend shall write each chunk to the
  temporary before its \texttt{write\_body} call returns, holding no body
  bytes beyond the chunk being written (\reqref{FRE-COMPAT.UNSIZED-01}).
  Only then does the valid length advance, so a partial-hit reader, which
  sees only bytes already in the file, is woken as soon as the chunk is
  there (\reqref{FDS-STORAGE.PARTIAL-01}).

  \textbf{Publishing.} WHEN a store completes, \texttt{finish} shall, in
  order: patch \texttt{body\_len} and \texttt{body\_crc}, recompute
  \texttt{header\_crc}, synchronise the temporary's data to stable storage
  (\reqref{ADR-PERSISTENCE.DURABILITY-01}) unless stores are unsynchronised
  (\reqref{FDS-PERSISTENCE.CRASHDISCARD-01}), and rename the temporary over the
  canonical name, dropping the suffix, while holding the key's publication
  lock (\reqref{FDS-STORAGE.METAUPDATE-01}). The shard directory is not
  synchronised afterwards. The record the rename replaces is held open
  across it and closed last, once the lock is released: closing it frees
  its storage, which can take a long time for a large record.
  \texttt{finish} shall report the record size --- the published
  file's full length --- as the created
  size,\footnote{\pgsrc{pingora-cache/src/storage.rs}{363}{\texttt{HandleMiss::finish}}
  returns a \pgsrc{pingora-cache/src/storage.rs}{346}{\texttt{MissFinishType}}.}
  which \texttt{HttpCache} admits to the eviction manager
  (\reqref{FRE-COMPAT.EVICTION-01}). The rename is the object's \emph{durable
  publication} --- the point at which it becomes visible by name to a scan, a
  restart, or a peer process --- not the first moment it can be read, and an
  incomplete store therefore never becomes visible under the canonical name
  (\reqref{FRE-COMPLIANCE.INCOMPLETE-01}). Once the backend is read-only
  (\reqref{FDS-COMPAT.READONLY-01}), \texttt{finish} instead removes the
  temporary rather than renaming it, while still reporting success, so the
  draining process's request completes normally. WHEN an invalidating purge
  has ended the store (\reqref{FDS-STORAGE.PURGE-01}), \texttt{finish}
  shall likewise remove the temporary rather than rename it, deciding under
  the publication lock the purge marked the store under, and shall report
  an error, so that \texttt{HttpCache} admits nothing to its eviction
  manager and lets the requests waiting on its cache lock
  retry.\footnote{\pgsrc{pingora-cache/src/lib.rs}{1097}{\texttt{HttpCache::finish\_miss\_handler}}
  releases its cache lock as a transient error when \texttt{finish} fails,
  and returns before admitting the entry
  (\pgsrc{pingora-cache/src/lib.rs}{1105}{L1105}).}

  \textbf{Failure and abort.} IF a write, the synchronisation, or the rename
  fails during a store, THEN the \texttt{write\_body} or \texttt{finish} call
  in progress shall return that error (\reqref{FRE-COMPAT.WRITEFAIL-01}), and
  the store is treated as aborted from then on. WHEN a store ends without
  being published --- after such an error, after an upstream failure or
  cancellation in which \texttt{HttpCache} drops the miss handler, once the
  backend is read-only, or once a purge has ended the store --- the backend
  shall unlink the temporary no later than when the
  miss handler is dropped (\reqref{FRE-COMPAT.ABORTSPACE-01}). Its own handle
  on the temporary stays open across the unlink, so the unlink removes only
  the name, and it closes that handle on the blocking pool
  (\reqref{FDS-STORAGE.EXECUTOR-01}): the close frees the storage, which
  can take a long time for a large temporary. A partial-hit
  reader still holding the temporary open is unaffected by the unlink; how its
  stream ends is specified by \reqref{FDS-STORAGE.PARTIAL-01}.

  \textbf{Isolation and replacement.} Same-key writers stay isolated under
  this scheme: each has its own temporary, so no writer's bytes are ever
  interleaved with another's (\reqref{FRE-COMPAT.CONCURRENT-01}), and a
  temporary abandoned by a crash is identifiable by its extension, never
  served or promoted, and is reclaimed by the restart-time scan
  (\reqref{FDS-STORAGE.RECLAIM-01}, \reqref{FRE-COMPAT.STARTRECLAIM-01}). A
  replacement store (\reqref{FRE-COMPLIANCE.REPLACE-01}) is just another
  per-writer temporary racing to publish under the same canonical name;
  convergence on a single winner falls out of last-rename-wins, and a reader
  that already opened the outgoing object keeps reading it to completion
  undisturbed, since a rename never invalidates a file descriptor already
  open on the old name. A peer process observes an object only at its
  durable publication.

  Figure~\ref{fig:fds-publish-seq} drills into the \texttt{get\_miss\_handler}
  / \texttt{write\_body} / \texttt{finish} calls that Figure~\ref{fig:cache-miss-seq}
  shows as opaque calls to \texttt{Storage}, detailing what this backend does
  inside each one.

  \begin{figure}[H]
    \centering
    \begin{adjustbox}{max width=\textwidth}
    \begin{sequencediagram}
      \newinst[0.0]{HttpCache}{HttpCache}
      \newinst[3.2]{Storage}{Storage}
      \newinst[5.1]{Follower}{Follower}
      \begin{call}{HttpCache}{get\_miss\_handler(key, meta)}{Storage}{MissHandler}
        \postlevel
        \postlevel
        \postlevel
        \begin{callself}{Storage}{{\small \shortstack{create per-writer temporary exclusively;\\write header + key + metadata region;\\register the in-flight write}}}{}
        \end{callself}
      \end{call}
      \begin{sdalt}{loop}{[for each chunk]}
        \begin{call}{HttpCache}{write\_body(chunk, eof)}{Storage}{}
          \postlevel
          \postlevel
          \begin{callself}{Storage}{{\small \shortstack{append to temporary; update body CRC;\\advance valid length; wake readers}}}{}
          \end{callself}
        \end{call}
        \begin{call}{Follower}{read\_body() (concurrent partial hit)}{Storage}{chunk}
        \end{call}
      \end{sdalt}
      \begin{call}{HttpCache}{finish()}{Storage}{record size}
        \postlevel
        \postlevel
        \postlevel
        \begin{callself}{Storage}{{\small \shortstack{patch body\_len, body\_crc, header\_crc;\\fdatasync temporary;\\rename over canonical name}}}{}
        \end{callself}
      \end{call}
    \end{sequencediagram}
    \end{adjustbox}
    \caption{Publication sequence, drilling into the calls
    Figure~\ref{fig:cache-miss-seq} shows opaque: \texttt{get\_miss\_handler}
    creates the per-writer temporary, writes its header, key, and metadata
    region, and registers the write for partial hits
    (\reqref{FDS-STORAGE.PARTIAL-01}); each \texttt{write\_body} call appends
    a chunk and advances the valid length a concurrent \texttt{Follower}
    streams from; and \texttt{finish} patches the header, synchronises the
    temporary unless stores are unsynchronised, and renames it over the
    canonical name --- the object's
    durable publication --- returning the record size.}\label{fig:fds-publish-seq}
  \end{figure}

  \reqdiaglink{overview}{\textit{View the domain overview.}}

  \reqcover{verify}{TC}
  \reqcover{satisfy}{IMP}

  \refine{ADR-STORAGE.OBJECTFILE-01}
  \refine{ADR-STORAGE.ATOMICPUBLISH-01}
  \refine{ADR-STORAGE.WRITERID-01}
  \refine{ADR-PERSISTENCE.DURABILITY-01}
  \satisfy{FRE-COMPAT.CONCURRENT-01}
  \satisfy{FRE-COMPLIANCE.INCOMPLETE-01}
  \satisfy{FRE-COMPLIANCE.REPLACE-01}
  \satisfy{FRE-COMPLIANCE.READERISOLATION-01}
  \satisfy{FRE-COMPAT.WRITEFAIL-01}
  \satisfy{FRE-COMPAT.ABORTSPACE-01}
  \satisfy{FRE-COMPAT.UNSIZED-01}
  \trace{FDS-STORAGE.PARTIAL-01}
  \trace{SDS-LAYOUT.SHARDING-01}
\end{requirement}

\begin{requirement}{FDS-STORAGE.PARTIAL-01}[draft][functionalRequirement]{Streaming partial hits}

  \localreqdiagram*{fdspartial01}{\reqid}{\reqdesc}

  \textbf{In-flight registry.} The file system does not signal appends, so
  partial hits need coordination beside the file. The backend shall keep, in
  process memory, a registry of the stores in flight, keyed by the object's
  combined hash; for each store it holds the temporary's path, the object's
  \texttt{CacheMeta}, the valid body length written so far, a notification
  channel, and the order in which the store opened. A store enters the
  registry when \texttt{get\_miss\_handler} opens its temporary
  (\reqref{FDS-STORAGE.PUBLISH-01}), and
  leaves it when the store completes
  --- published, or removed once the backend is read-only or an invalidating
  purge has ended it --- or aborts. An invalidating purge marks the key's
  stores in flight as ended (\reqref{FDS-STORAGE.PURGE-01}); they stay in the
  registry until they complete or abort. The
  backend's \texttt{support\_streaming\_partial\_write}
  shall return \texttt{true} (\reqref{FRE-COMPAT.STREAMFLAG-01}).

  \textbf{Lookup order.} WHEN \texttt{lookup} is called, the backend shall
  first consult the registry. If one or more stores of the key that no
  invalidating purge has ended are in flight, it shall return a partial hit
  backed by the most recently opened of them; otherwise it shall look up the
  published record as \reqref{SDS-LAYOUT.OBJECTFILE-01} specifies. The
  request writing a store reads its body back through a lookup that names
  the store by its streaming write
  tag;\footnote{\pgsrc{pingora-cache/src/storage.rs}{166}{\texttt{Storage::lookup\_streaming\_write}},
  which \pgsrc{pingora-cache/src/lib.rs}{1031}{\texttt{HttpCache::set\_miss\_handler}}
  calls with the tag of the store it has just opened.} that lookup shall be
  answered with the store it names, or with nothing, and with the store even
  once a purge has ended it.

  \textbf{Partial-hit readers.} A partial hit's reader opens its own
  descriptor on the store's temporary and delivers the bytes up to the valid
  length, then waits on the store's notification channel. WHEN
  \texttt{write\_body} appends a chunk, the backend shall advance the valid
  length and wake every reader of that store before \texttt{write\_body}
  returns, so that each reader can deliver the chunk within the latency bound
  \reqref{FRE-COMPAT.STREAMDELIVERY-01} sets. WHEN the store completes ---
  whether it is published, or removed once the backend is read-only or a
  purge has ended it --- each
  reader shall
  deliver the remaining bytes and end normally
  (\reqref{FRE-COMPAT.STREAMEND-01}); its descriptor stays valid across the
  rename and across the removal. WHEN the store aborts, each reader shall end
  its stream with an error, never a clean end of body
  (\reqref{FRE-COMPLIANCE.PARTIALABORT-01}). A partial hit does not support
  seeking and, like the stock backend's, counts no eviction
  access.\footnote{\pgsrc{pingora-cache/src/memory.rs}{211}{\texttt{should\_count\_access}}
  returns \texttt{false} for the stock backend's partial reads.}

  \textbf{Scope.} The registry belongs to one process. A peer process ---
  such as the other side of a hot upgrade --- observes an object only at its
  durable publication (\reqref{ADR-STORAGE.ATOMICPUBLISH-01}).

  \reqdiaglink{overview}{\textit{View the domain overview.}}

  \reqcover{verify}{TC}
  \reqcover{satisfy}{IMP}

  \refine{ADR-STORAGE.OBJECTFILE-01}
  \satisfy{FRE-COMPAT.STREAMFLAG-01}
  \satisfy{FRE-COMPAT.STREAMING-01}
  \satisfy{FRE-COMPAT.STREAMDELIVERY-01}
  \satisfy{FRE-COMPAT.STREAMEND-01}
  \satisfy{FRE-COMPLIANCE.PARTIALABORT-01}
  \trace{FDS-STORAGE.PUBLISH-01}
\end{requirement}

\begin{requirement}{FDS-STORAGE.METAUPDATE-01}[draft][functionalRequirement]{Metadata update: in place or withdrawn}

  \localreqdiagram*{fdsmetaupdate01}{\reqid}{\reqdesc}

  Each object's metadata sits in a region of its record, right after the key
  and ahead of the body (\reqref{ADR-STORAGE.OBJECTFILE-01}), whose capacity
  the record's header holds (\reqref{SDS-LAYOUT.OBJECTFILE-01}). Unlike the body, whose length is
  unknown until the stream ends (\reqref{FRE-COMPAT.UNSIZED-01}), the
  metadata's exact size is already known when the record is created ---
  \texttt{CacheMeta} is fully serialised at that point --- so the region
  reserves exactly that size plus a \emph{configurable} amount of growth
  headroom beyond it --- 512 bytes by default
  (\reqref{SDS-LAYOUT.OBJECTFILE-01}) --- to anticipate a later metadata
  update. How much
  headroom to reserve is deployment-tunable, trading reserved-but-possibly-unused
  space against how often a later update fits in place; a
  deployment configured with zero headroom simply withdraws the record, as
  described below, on every metadata change that grows it, rather than
  needing a separate mechanism for it.

  A metadata update (\reqref{FRE-COMPLIANCE.REVALIDATION-01}) whose new
  segments still fit the region shall be applied
  in place: the segments are rewritten within the region and the record's recorded
  lengths and checksums updated, without touching the body and without
  synchronising to stable storage (\reqref{ADR-PERSISTENCE.DURABILITY-01}).

  The backend shall open and read the stored record, decide which of the two
  paths applies, and write the new metadata in place or withdraw the record,
  as described below, all while holding the key's
  \emph{publication lock}. Every change that replaces
  or removes the key's record --- a store's rename, a purge's
  or the reclaim scan's unlink --- holds that lock too, so the update works on
  the record the key holds once the lock is taken. A record opened before the
  lock could lose its name in between, and the update would change a file no
  lookup reads any more, yet report success. Reading outside that lock would
  let one update observe another's region write half-done. An update that
  read a torn region could not distinguish a racing writer from a damaged
  record, and would report a perfectly healthy object as absent
  (\reqref{FRE-COMPLIANCE.METAABSENT-01}) rather than updating it. A
  \emph{reader} re-reads instead, and takes the lock only for a last read of
  a record whose checksums kept failing, which then finds it whole
  (\reqref{SDS-LAYOUT.OBJECTFILE-01}).

  Before either path, still holding the lock, the backend shall check that
  the record holds the response the update freshens. \texttt{pingora-cache}
  gives the new metadata the creation time of the response its \texttt{304}
  validated,\footnote{\pgsrc{pingora-cache/src/lib.rs}{1200}{\texttt{HttpCache::revalidate\_cache\_meta}}
  copies the stale response's \texttt{created} into the new metadata.} so a
  stored response created at another time has replaced that one since the
  lookup that began the revalidation. \texttt{pingora-cache} permits that
  once a writer outlives its cache lock,\footnote{\pgsrc{pingora-cache/src/lock.rs}{110}{\texttt{CacheLock::lock}}
  hands a new write permit to a waiter once the holder's lock has passed its
  age timeout.} or when no cache lock is configured. The update shall then not
  be applied: \texttt{update\_meta} reports \texttt{false}, and the newer
  response stands whole (\reqref{FRE-COMPLIANCE.REPLACE-01}). Applied, it
  would pair the validated response's header fields with the newer response's
  body, which \rfcsec{9111}{4.3.4} rules out: a \texttt{304} whose strong
  validators no stored response shares must not update any.

  When an update's new segments no longer fit the region's capacity ---
  whether because the growth is unusually large or because no headroom was
  configured --- the backend shall withdraw the record instead, so that the
  next request fetches the object again: it rewrites only the record's
  header, marked withdrawn (\reqref{SDS-LAYOUT.OBJECTFILE-01}), and reports
  the update not applied. Every lookup then misses, and Pingora fetches the
  object from the origin and stores it, through a store whose region is
  sized for the new metadata and whose rename replaces the file
  (\reqref{FDS-STORAGE.PUBLISH-01}). The response that triggered the update
  is still served, from the record its hit handler holds open. Until the
  store replaces it, the withdrawn file keeps its length, so the eviction
  manager, which the backend never calls (\reqref{FRE-COMPAT.EVICTION-01}),
  keeps counting the space it takes; an eviction or a purge removes it as it
  removes any record, and the restart procedure reclaims it
  (\reqref{FDS-STORAGE.RECLAIM-01}). The body is neither read nor moved. A
  withdrawal whose header write fails reports an error and leaves the
  record as it was. A withdrawal is not synchronised to stable storage
  either (\reqref{ADR-PERSISTENCE.DURABILITY-01}): after a host crash the
  record is withdrawn, as it was, or torn, which a reader rejects. Moving
  outgrown metadata after the body instead, which would spare the fetch, is
  deferred to a later version (\reqref{ADR-STORAGE.METAMOVE-01}).

  \texttt{update\_meta} reports \texttt{true} once the update is written in
  place, and so visible to lookups, and \texttt{false} when it is not
  applied: a newer response has replaced the one it validated, or the
  record was withdrawn.
  The publication lock is a per-key lock
  in process memory, which a store's final rename
  (\reqref{FDS-STORAGE.PUBLISH-01}), a purge's unlink
  (\reqref{FDS-STORAGE.PURGE-01}), the reclaim scan's unlink of a canonical
  record (\reqref{FDS-STORAGE.RECLAIM-01}), and the restart procedure's
  admission checks (\reqref{FDS-EVICTION.STARTUP-01}) also take, so none can
  land between a confirmation and the action that depends on it; a lookup
  takes it only for the last read described above. Across processes, only
  the hot-upgrade window before the outgoing backend becomes read-only
  (\reqref{FDS-COMPAT.READONLY-01}) lets two backends write, and the residual
  race there is accepted. Once the backend is read-only, \texttt{update\_meta}
  fails with an error on either path (\reqref{FRE-COMPAT.HOTUPGRADE-01}).
  With a non-zero configured headroom,
  the common freshening after a \texttt{304} (\texttt{Date}- and
  \texttt{Age}-scale changes) is small enough to take the in-place path;
  growth beyond whatever headroom is configured always withdraws the
  record.

  A stale hit's revalidation succeeds when the origin returns a
  \texttt{304} (\reqref{FRE-COMPLIANCE.REVALIDATION-01}).
  \texttt{HttpProxy}'s \texttt{revalidate\_or\_stale} merges the
  \texttt{304} into the cached header and calls into \texttt{HttpCache}'s
  metadata-update handling, which resolves to one of the two execution
  paths above.

  When a \texttt{CacheLock} is configured (it is an optional collaborator),
  \texttt{HttpCache} already holds the key's write lock by this point:
  \texttt{cache\_lookup} takes it as soon as the hit is first found stale,
  well before the upstream request even starts, so that other requests
  for the same key queue behind this one leader. The lock is released
  only after the metadata-update call returns, whether it succeeds or
  fails.\footnote{\pgsrc{pingora-cache/src/lib.rs}{1539}{\texttt{HttpCache::cache\_lookup}}
  takes the lock when storage returns nothing; the stale path does so in
  \pgsrc{pingora-cache/src/lib.rs}{862}{\texttt{cache\_found}}.
  \pgsrc{pingora-proxy/src/proxy_cache.rs}{774}{\texttt{HttpProxy::revalidate\_or\_stale}}
  merges the response
  (\pgsrc{pingora-proxy/src/proxy_cache.rs}{808}{\texttt{revalidate\_merge\_header}})
  and calls
  \pgsrc{pingora-proxy/src/proxy_cache.rs}{825}{\texttt{revalidate\_cache\_meta}};
  \pgsrc{pingora-cache/src/lib.rs}{1187}{\texttt{HttpCache::revalidate\_cache\_meta}}
  calls \pgsrc{pingora-cache/src/lib.rs}{1218}{\texttt{update\_meta}} and
  releases the lock regardless of the result
  (\pgsrc{pingora-cache/src/lib.rs}{1230}{L1230}).}

  \texttt{update\_meta} cannot manipulate the lock.\footnote{\pgsrc{pingora-cache/src/storage.rs}{235}{Its signature}
  is \texttt{(key, meta, trace) -> Result<bool>}: no reference to the
  lock is passed in.} It shall make the update, or the withdrawal, visible
  to lookups before returning. \texttt{HttpCache} releases the lock as soon
  as this call returns (Figure~\ref{fig:fds-metaupdate-seq}), so every
  lookup after it, the followers queued behind the lock included, must find
  the update, or miss. On either path the change is complete when it is
  visible.

  Lookups do not wait for the lock: a fresh hit never takes it, and a
  stale response may be served while a background subrequest revalidates
  it.\footnote{Under stale-while-revalidate, the lock writer in
  \pgsrc{pingora-proxy/src/proxy_cache.rs}{243}{\texttt{HttpProxy::proxy\_cache}}
  hands its write lock to a background subrequest and serves the stale
  object, and a request finding the key locked serves stale too
  (\pgsrc{pingora-proxy/src/proxy_cache.rs}{233}{L233}) instead of
  waiting.} A
  concurrent lookup can therefore read the object while either path runs.
  The in-place path writes into the region a lookup may be reading, and a
  withdrawal into the header; a lookup that catches either mid-write is
  protected by the read-side validation and retry that
  \reqref{SDS-LAYOUT.OBJECTFILE-01} specifies, so it observes the
  pre-update metadata, the post-update metadata, or a miss once the record
  is withdrawn, never a mixture. The body is untouched on either path.

  A withdrawal costs one header write, whatever the body's size, and then
  one fetch of the whole object from the origin. A herd of revalidations
  withdraws the record once: the first update that outgrows the region
  withdraws it, and the next ones find no object stored, which they report
  as an update of an absent key does (\reqref{FRE-COMPLIANCE.METAABSENT-01}).
  The configurable headroom above keeps withdrawals rare, and they are
  observable, as \reqref{FDS-OBSERVABILITY.HEADROOM-01} specifies.

  Figure~\ref{fig:fds-metaupdate-seq} shows the whole sequence, using the
  same \texttt{HttpProxy}, \texttt{HttpCache}, \texttt{CacheLock}, and
  \texttt{Storage} structure as the cache-miss sequence in the System
  Context chapter (Figure~\ref{fig:cache-miss-seq}).

  \begin{figure}[H]
    \centering
    \begin{adjustbox}{max width=\textwidth}
    \begin{sequencediagram}
      \newinst[0.0]{HttpProxy}{HttpProxy}
      \newinst[5.2]{HttpCache}{HttpCache}
      \newinst[2.8]{CacheLock}{CacheLock}
      \newinst[0.5]{Storage}{Storage}
      \begin{call}{HttpProxy}{cache\_lookup()}{HttpCache}{stale hit}
        \begin{call}{HttpCache}{lookup(key)}{Storage}{stale}
        \end{call}
        \begin{call}{HttpCache}{lock(key, stale)}{CacheLock}{leader (write permit)}
        \end{call}
      \end{call}
      \postlevel
      \begin{callself}{HttpProxy}{{\small fetch upstream; receives 304 (Not Modified)}}{}
      \end{callself}
      \postlevel
      \begin{callself}{HttpProxy}{revalidate\_or\_stale(task)}{}
        \begin{call}{HttpProxy}{revalidate\_merge\_header(resp)}{HttpCache}{merged headers}
        \end{call}
        \begin{call}{HttpProxy}{revalidate\_cache\_meta(new\_meta)}{HttpCache}{}
          \begin{sdalt}{alt}{[new segments fit the metadata region]}
            \begin{call}{HttpCache}{update\_meta(key, meta)}{Storage}{true: updated in place}
            \end{call}
          \sdpart{[new segments do not fit]}
            \begin{call}{HttpCache}{update\_meta(key, meta)}{Storage}{false: withdrawn}
              \postlevel
              \postlevel
              \postlevel
              \begin{callself}{Storage}{{\small \shortstack{under the publication lock: rewrite the header,\\marked withdrawn; the next lookup misses,\\and the object is fetched again}}}{}
              \end{callself}
            \end{call}
          \end{sdalt}
          \begin{call}{HttpCache}{release(key, permit, Done)}{CacheLock}{}
          \end{call}
        \end{call}
      \end{callself}
    \end{sequencediagram}
    \end{adjustbox}
    \caption{Metadata-update sequence, in the same \texttt{HttpProxy},
    \texttt{HttpCache}, \texttt{CacheLock}, and \texttt{Storage}
    structure as Figure~\ref{fig:cache-miss-seq}. The write lock is
    already held when \texttt{revalidate\_or\_stale} begins and is
    released only after \texttt{update\_meta} returns, whichever path it
    took. On either path a header write, made last, commits the change, and
    the body stays where it is.}\label{fig:fds-metaupdate-seq}
  \end{figure}

  \reqdiaglink{overview}{\textit{View the domain overview.}}

  \reqexpect{verify}{TC-COMPLIANCE.REVALIDATION-01}{Extend this TC to also
  verify FDS-STORAGE.METAUPDATE-01: that an update written in place is
  visible to lookups when update_meta returns; that metadata outgrowing its
  region withdraws the record, which lookups then miss while a client
  already reading it gets its whole body; that the next store replaces a
  withdrawn record; that a withdrawal whose header write fails leaves the
  record as it was; that an update applies only to the response it
  validated; and that zero configured headroom withdraws the record on any
  growth.}
  \reqcover{satisfy}{IMP}

  \refine{ADR-STORAGE.OBJECTFILE-01}
  \refine{ADR-PERSISTENCE.DURABILITY-01}
  \satisfy{FRE-COMPLIANCE.REVALIDATION-01}
  \satisfy{FRE-COMPLIANCE.METAATOMIC-01}
  \satisfy{FRE-COMPLIANCE.METAPERSIST-01}
  \satisfy{FRE-COMPLIANCE.METAABSENT-01}
  \trace{FDS-STORAGE.PUBLISH-01}
  \trace{FDS-STORAGE.PURGE-01}
  \trace{REQ-EFFICIENCY.FRAGMENTATION-01}
\end{requirement}

\begin{requirement}{FDS-OBSERVABILITY.HEADROOM-01}[draft][functionalRequirement]{Headroom-exceeded observability}

  \localreqdiagram*{fdsobsheadroom01}{\reqid}{\reqdesc}

  WHEN a metadata update finds that its new metadata outgrew the record's
  reserved headroom, so that the update withdraws the record
  (\reqref{FDS-STORAGE.METAUPDATE-01}), the backend shall log a warning and, when
  its \texttt{metrics} feature is enabled (it is off by default,
  \reqref{ADR-OBSERVABILITY.METRICS-01}), increase a counter:
  \begin{itemize}
    \item The warning, emitted through the \texttt{log} facade
      \texttt{pingora-cache} already depends
      on,\footnote{\pgsrc{pingora-cache/Cargo.toml}{37}{\texttt{log}} is a
      dependency of \texttt{pingora-cache}; its own
      \pgsrc{pingora-proxy/src/proxy_cache.rs}{828}{\texttt{warn!("revalidate\_cache\_meta failed...")}}
      is precedent for logging a metadata-update outcome.} names the record
      by its path, whose file name is the key's hex-encoded combined
      hash\footnote{\pgsrc{pingora-cache/src/key.rs}{90}{\texttt{CacheHashKey::combined}}:
      the only human-inspectable identity \texttt{Storage::update\_meta}'s
      \texttt{CacheKey} exposes is a hash, not the original request URL.}
      and states the headroom shortfall: how many more bytes of headroom
      would have kept the update in place.
    \item The counter \texttt{pingora\_durable\_meta\_headroom\_exceeded\_total},
      registered in the \texttt{prometheus} default registry
      (\reqref{ADR-OBSERVABILITY.METRICS-01}) without any per-key label, so
      how often it happens is visible without per-key cardinality, is
      increased by one. The backend registers it, at zero, when it is
      built: a counter that appeared only with its first event could not
      be told from one that is not served at all, and \texttt{rate()} and
      \texttt{increase()} would miss the event that created it.
  \end{itemize}
  Both are emitted once the record is withdrawn, so a withdrawal
  whose header write fails, and so did not happen, is neither logged nor
  counted. With the \texttt{metrics} feature disabled, the counter is not
  registered and the warning is the only record of each case. A count that
  keeps rising tells an operator to raise the configured headroom.

  \reqdiaglink{overview}{\textit{View the domain overview.}}

  \reqcover{verify}{TC}
  \reqcover{satisfy}{IMP}

  \refine{ADR-OBSERVABILITY.METRICS-01}
  \refine{FDS-STORAGE.METAUPDATE-01}
\end{requirement}

\begin{requirement}{ADR-STORAGE.METAMOVE-01}[draft][designConstraint]{Metadata moves deferred for v1}

  \localreqdiagram*{adrmetamove01}{\reqid}{\reqdesc}

  \textbf{Context:} A metadata update whose new metadata outgrows the region
  reserved ahead of the body cannot be written in place
  (\reqref{FDS-STORAGE.METAUPDATE-01}). The backend then withdraws the
  record: the response at hand is still served, the next request fetches
  the object from the origin, and its store, sized for the new metadata,
  replaces the file. That costs a fetch of the whole body, and, should the
  origin fail before it, the stale copy that could have been served
  instead. The configurable headroom keeps it rare, and every case is
  observable (\reqref{FDS-OBSERVABILITY.HEADROOM-01}), so an operator can
  raise the headroom where it is not.

  \textbf{Decision:} Do not move outgrown metadata for this version. The
  first alternative below would spare the fetch, but it changes the record
  format and the reader's rules, and it needs the backend to tell the
  eviction manager how much each move grows a record, which it cannot do
  today: the manager is the framework's, and the backend holds it only during
  its restart procedure (\reqref{FRE-COMPAT.EVICTION-01}).
  \reqref{HLR-SCOPE.V1-01} prefers the simplest adequate design, so the
  fetch is accepted for this version, and moving the metadata is to be
  revisited in a later one.

  \textbf{Alternatives considered:}
  \begin{itemize}
    \item \emph{Move the metadata after the body}: write the new metadata
      into a slot appended after the end of the record, with headroom, then
      the header that points to it. The body is never read or moved, and the
      old metadata stays where it was, so a lookup that read the old header
      reads it whole. It needs two more header fields, saying where the
      metadata and the body lie; a reader that accepts a file longer than
      its record, since a move writes its slot before its header; and a
      report to the eviction manager of each move's growth, through
      \texttt{EvictionManager::increment\_weight}, purging the victims it
      returns. Deferred: those changes, for a case the headroom makes rare.
    \item \emph{Store the object again from its own body}, through a store
      of the key opened with the new metadata and a larger region. Rejected:
      it copies a whole body to change its metadata, and lookups follow the
      copy while it runs, so a purge, an expiry, or a newer store that lands
      meanwhile is not seen until it ends.
    \item \emph{Withdraw the record} (chosen): as the context describes.
  \end{itemize}

  \textbf{Consequences:}
  \begin{itemize}
    \item[+] The record format keeps the metadata in one place, and a reader
      checks one exact length.
    \item[+] The eviction manager stays in step without the backend calling
      it: a withdrawn record keeps its file, and the space the manager
      counts for it, until a store replaces it.
    \item[--] An update that outgrows the headroom costs a miss for the next
      request and a fetch of the whole object from the origin.
    \item[--] Between the withdrawal and that fetch, no stale copy is left to
      serve should the origin fail.
  \end{itemize}

  \reqdiaglink{overview}{\textit{View the domain overview.}}

  \reqterminal[Deferred: metadata moves are not implemented for this version,
  so nothing further refines, satisfies, or verifies this decision; revisit
  in a later version.]

  \refine{ADR-STORAGE.OBJECTFILE-01}
  \trace{HLR-SCOPE.V1-01}
  \trace{FDS-STORAGE.METAUPDATE-01}
  \trace{FRE-COMPAT.EVICTION-01}
\end{requirement}

\begin{requirement}{ADR-OBSERVABILITY.METRICS-01}[draft][designConstraint]{Metrics via the prometheus crate}

  \localreqdiagram*{adrmetrics01}{\reqid}{\reqdesc}

  \textbf{Context:} The backend accepts some costs rather than engineering
  them away but makes them observable --- for example the fetch from the
  origin that a metadata update causes when it outgrows its reserved
  headroom (\reqref{FDS-STORAGE.METAUPDATE-01}). Aggregate counters need a metrics
  library, and the choice decides where an operator sees them. Pingora ships an
  HTTP application that exports the \texttt{prometheus} crate's default
  registry,\footnote{\pgsrc{pingora-prometheus/src/lib.rs}{82}{\texttt{PrometheusHttpApp}}
  serves \pgsrc{pingora-prometheus/src/lib.rs}{88}{\texttt{prometheus::gather()}},
  the crate's default registry.} so registering there is what puts the backend's
  counters in front of an operator who already scrapes Pingora.

  Where that application lives changed in Pingora 0.9.0, and with it the cost of
  this decision. Under 0.8.1 the exporter was the
  \texttt{prometheus\_http\_app} module of \texttt{pingora-core} itself, and \texttt{prometheus} was an
  unconditional dependency of \texttt{pingora-core}, so every build that used
  \texttt{pingora-cache}\footnote{\pgsrc{pingora-cache/Cargo.toml}{21}{\texttt{pingora-core}}
  is an unconditional dependency of \texttt{pingora-cache}.} already contained
  the crate and no feature of this backend's could remove it. In 0.9.0 the
  exporter moved to a separate \texttt{pingora-prometheus} crate and
  \texttt{prometheus} is no longer a \texttt{pingora-core} dependency at all.
  A build of this backend that registers metrics therefore now pulls
  \texttt{prometheus} in on its own account. The exporter crate re-exports
  \texttt{prometheus} for exactly this use,\footnote{\pgsrc{pingora-prometheus/src/lib.rs}{54}{``Re-export
  of the \texttt{prometheus} crate''}: ``Use this re-export to ensure your
  metrics are registered in the same'' registry the exporter serves;
  \pgsrc{pingora-prometheus/src/lib.rs}{70}{\texttt{pub use prometheus}}.}
  and no Pingora crate has a feature flag for metrics that the backend's
  own choice could follow.

  \textbf{Decision:} Register the backend's metrics with the
  \texttt{prometheus} crate's default registry, through the
  \texttt{prometheus} crate that \texttt{pingora-prometheus} re-exports, so
  they appear wherever an integrator exposes the Pingora metrics --- but
  only when the backend's \texttt{metrics} Cargo feature is enabled, which it
  is not by default. The feature gates both the registration and the
  dependency: with it disabled the backend neither registers metrics nor
  links \texttt{pingora-prometheus} or \texttt{prometheus} at all. Going
  through the exporter's re-export makes the registry the exporter's by
  construction: the default registry is a static of the \texttt{prometheus}
  crate, and the backend cannot resolve to a different version of it than
  the exporter does. Enabling the feature is the same decision as running
  that exporter, since counters nothing exports are of no use. With the
  feature disabled, the backend registers no metrics and its warning logs
  remain the only record of the events the metrics would have counted.

  \textbf{Alternatives considered:}
  \begin{itemize}
    \item \emph{Register unconditionally, with no feature}: the simplest
      option. Rejected: it leaves an integrator no way to keep the backend's
      names out of a registry it manages itself, and --- since Pingora 0.9.0
      made \texttt{prometheus} optional --- it would also force the dependency
      on builds that do not want it.
    \item \emph{Gate the metrics behind a feature enabled by default}: this
      decision's previous revision. Rejected: every integrator who does not
      run the \texttt{pingora-prometheus} exporter would carry a dependency and
      counters nothing reads, while the events the counters measure are
      logged regardless.
    \item \emph{Follow a Pingora feature flag}: not available --- none of
      the Pingora 0.9.0 crates has a metrics feature, and Cargo cannot enable
      a feature because the integrator depends on some other crate.
    \item \emph{Depend on \texttt{prometheus} directly}: the previous
      revision did. Rejected: it registers into the exporter's registry only
      while both crates happen to resolve to the same \texttt{prometheus}
      version, which nothing checks.
    \item \emph{The \texttt{metrics} facade crate}, leaving the exporter to
      the integrator: rejected --- metrics recorded through it do not reach
      the \texttt{prometheus} default registry that the Pingora exporter
      serves without an additional bridge, and it would replace one optional
      dependency with a different one plus that bridge.
    \item \emph{Logs only, no counters}: rejected --- a warning per event
      shows individual occurrences, but not the aggregate volume the
      design needs visible without per-key cardinality.
  \end{itemize}

  \textbf{Consequences:}
  \begin{itemize}
    \item[+] With the feature on, metrics appear alongside the Pingora
      metrics, in the registry the \texttt{pingora-prometheus} exporter
      serves, by construction.
    \item[+] By default the backend links neither \texttt{pingora-prometheus}
      nor \texttt{prometheus}, and registers nothing.
    \item[--] The counters exist only where an integrator enables the
      feature; elsewhere the costs they measure are visible in the warning
      logs alone.
    \item[--] With the feature on, the deliverable depends on
      \texttt{pingora-prometheus}, which brings the exporter's code along with
      the re-export, even for an integrator who serves the registry another
      way.
    \item[--] The metrics live in a process-global registry, so their names
      must not collide with those of the integrator or of Pingora.
    \item[--] Two build configurations to keep compiling and tested.
  \end{itemize}

  \reqdiaglink{overview}{\textit{View the domain overview.}}

  \reqexpect{refine}{FDS-OBSERVABILITY.HEADROOM-01}{Specify each metric
  emitted when new metadata outgrows its record's headroom --- name, type,
  unit, and when it changes --- the accompanying warning log, and what
  remains when the metrics feature is disabled.}

  \satisfy{HLR-SCOPE.V1-01}
  \trace{FDS-STORAGE.METAUPDATE-01}
\end{requirement}

\begin{requirement}{ADR-STORAGE.EXPIRE-01}[draft][designConstraint]{Expire by rewriting the metadata}

  \localreqdiagram*{adrexpire01}{\reqid}{\reqdesc}

  \textbf{Context:} \reqref{FRE-COMPAT.EXPIRE-01} asks that an expired object
  stop being served as fresh while its body stays available for a conditional
  revalidation. \texttt{Storage::expire} supplies a default that deletes the
  object instead, which meets the first half and discards the second.
  Pingora is explicit that how an expiry is recorded is the backend's to
  choose, and that rewriting the stored freshness and applying an expiry at
  read time both satisfy the trait's
  contract.\footnote{\pgsrc{pingora-cache/src/storage.rs}{226}{\texttt{Storage::expire}}:
  ``How that is recorded is up to storage: rewriting the stored freshness and
  applying it on read both satisfy the contract.''} The framework also supplies
  the transformation itself:
  \texttt{CacheMeta::expire\_at}\footnote{\pgsrc{pingora-cache/src/meta.rs}{800}{\texttt{CacheMeta::expire\_at}}.}
  produces the expired metadata from the stored metadata, without touching the
  body.

  What makes the choice cheap here is that this backend already has the
  machinery. An object's metadata sits in a region with configured growth
  headroom, and \reqref{FDS-STORAGE.METAUPDATE-01} already specifies how to
  replace it in place when the new segments fit the region, without reading
  or copying the body. An expiry's new segments differ from the stored ones
  only in the internal-metadata segment's freshness deadline --- the header
  segment is untouched --- yet they need not be the same length.
  \texttt{pingora-cache} serialises that segment with MessagePack, which
  writes an integer in as few bytes as its value
  allows,\footnote{One byte from 0 to 127, and two, three, five or nine above
  (\href{https://github.com/msgpack/msgpack/blob/master/spec.md\#int-format-family}{MessagePack
  specification: int format family}); the serialiser \texttt{pingora-cache}
  uses picks the shortest form that holds the value.} and a deadline's
  nanoseconds are one such integer. A deadline taken from an
  \texttt{Expires} header falls on a whole
  second,\footnote{\pgsrc{pingora-cache/src/filters.rs}{120}{\texttt{calculate\_expires\_header\_time}}
  returns the header's date, which has no fraction of a second.} so its
  nanoseconds take one byte, where those of an arbitrary instant nearly
  always take five: expired at the current instant, such an object's
  metadata grows by four bytes.

  \textbf{Decision:} Implement \texttt{expire} natively, by deriving the
  expired metadata with \texttt{expire\_at} at the start of the current
  second and writing it in place, as \reqref{FDS-STORAGE.METAUPDATE-01}
  writes an update that fits, rather than inheriting the deleting default.
  \texttt{expire\_at} only ever lowers the deadline, so its seconds take no
  more bytes than before, and a deadline on a whole second has nanoseconds
  of zero, which take the fewest: the expired metadata is never longer than
  the stored metadata, so it fits the region whatever headroom is left. The
  stored body is left untouched, so the \texttt{304} that follows reuses it.

  \textbf{Alternatives considered:}
  \begin{itemize}
    \item \emph{Inherit the default and delete}: the cheapest to specify and to
      build. Rejected --- it turns every expiry into a full refetch of a body
      the cache is still holding, which is the cost an expiring purge exists to
      avoid; a durable backend, whose whole argument is that bodies are large
      and worth keeping (\reqref{HLR-PERSISTENCE.DURABLE-01}), is the last place
      to discard one.
    \item \emph{Apply the expiry at read time from a side marker} --- a
      tombstone file, or a flag in the record header that \texttt{lookup}
      consults: rejected --- it adds a second on-disk thing per key for the
      restart scan to classify and reclaim (\reqref{FDS-STORAGE.RECLAIM-01}),
      and a header flag would have to be reconciled against the freshness the
      metadata segment itself carries, leaving two sources of truth for one
      question.
    \item \emph{Expire at the current instant, and fall back when the expired
      metadata does not fit the region}: rejected --- the fallback costs more
      than the four bytes it is for. Withdrawing the record, as a metadata
      update does (\reqref{FDS-STORAGE.METAUPDATE-01}), would report the
      entry still stored while no lookup could reuse its body; removing it,
      as the default does, reports a purge, which must then be as durable as
      one (\reqref{FRE-COMPLIANCE.PURGEDURABLE-01}).
    \item \emph{Keep four bytes free in every region for the expiry}: a store
      adds them, and an update in place leaves them free. Rejected --- three
      writers would keep an invariant that the start of the second keeps
      within the expiry alone.
  \end{itemize}

  \textbf{Consequences:}
  \begin{itemize}
    \item[+] An expiry costs one metadata region's worth of writes, never a
      body copy, and the retained body answers the revalidating \texttt{304}.
    \item[+] No new on-disk structure, no new reclaim class, and no new lock:
      the operation reuses the publication lock and the in-place path
      \reqref{FDS-STORAGE.METAUPDATE-01} already specifies.
    \item[--] An expiry inherits that path's crash semantics: after a host crash
      an in-place expiry reads back as the old metadata, the new metadata, or a
      miss (\reqref{ADR-PERSISTENCE.DURABILITY-01}), so an acknowledged expiry
      can be lost. This is weaker than the purge-acknowledgement guarantee
      \reqref{FRE-COMPLIANCE.PURGEDURABLE-01} makes, and is why an expiry is not
      a substitute for a purge when the object must provably stop being served.
    \item[+] An expiry never removes an object, whatever headroom its region
      has left.
    \item[--] An expired object is stale from the start of the second in
      which it was expired, so its serve-stale windows, measured from its
      deadline, close up to a second earlier than they would from the
      instant itself. HTTP caching works in whole
      seconds,\footnote{\pgsrc{pingora-cache/src/meta.rs}{716}{\texttt{CacheMeta::is\_fresh}}:
      ``HTTP cache time resolution is second''.} and Pingora itself backdates
      the deadline of a response with no freshness lifetime by a
      second.\footnote{\pgsrc{pingora-cache/src/filters.rs}{79}{\texttt{freshness\_ttl\_to\_time}}.}
  \end{itemize}

  \reqdiaglink{overview}{\textit{View the domain overview.}}

  \reqexpect{refine}{FDS-STORAGE.EXPIRE-01}{Specify how a target is resolved,
  how the expired metadata is derived and written in place, why it always
  fits, which outcome each case reports, and how the operation behaves once
  the backend is read-only.}

  \satisfy{FRE-COMPAT.EXPIRE-01}
  \refine{ADR-STORAGE.OBJECTFILE-01}
  \trace{FDS-STORAGE.METAUPDATE-01}
  \trace{FDS-STORAGE.PURGE-01}
  \trace{ADR-PERSISTENCE.DURABILITY-01}
\end{requirement}

\begin{requirement}{FDS-STORAGE.PURGE-01}[draft][functionalRequirement]{Purge via canonical-name unlink}

  \localreqdiagram*{fdspurge01}{\reqid}{\reqdesc}

  A purge is keyed directly off the presented \texttt{PurgeTarget}, with
  no lookup and no \texttt{CacheLock} (Figure~\ref{fig:purge-seq}). The target names the
  entry to remove in one of two ways: \texttt{Active}, carrying a
  \texttt{CompactCacheKey}, asks for whichever entry is currently active for
  that logical key --- how \texttt{HttpCache} issues an invalidation ---
  while \texttt{Exact}, carrying a \texttt{CacheEntryKey}, names one entry
  completely, which is how eviction victims arrive.\footnote{\pgsrc{pingora-cache/src/storage.rs}{51}{\texttt{PurgeTarget}};
  \pgsrc{pingora-cache/src/lib.rs}{1883}{\texttt{purge\_impl}} passes
  \texttt{Active} and \pgsrc{pingora-cache/src/lib.rs}{1130}{the eviction loop}
  passes \texttt{Exact}.} Because this backend assigns no entry identifiers
  (\reqref{ADR-COMPAT.ENTRYID-01}), the two carry the same information whenever
  the backend can act at all, and it shall resolve either --- an \texttt{Active}
  target, or an \texttt{Exact} target whose \texttt{CacheEntryKey} is
  \texttt{KeyOnly} --- to the canonical path of exactly that key, computed as
  \reqref{SDS-LAYOUT.SHARDING-01} specifies, from the hex encoding of the
  key's combined hash, just as the stock backend
  does\footnote{\pgsrc{pingora-cache/src/memory.rs}{402}{\texttt{MemCache::purge}}
  resolves the entry by \texttt{target.key().combined()}.} --- and shall never touch
  any other file. An \texttt{Exact} target carrying an \texttt{Identified}
  entry key names a generation this backend never created, so the backend
  shall report it as not found and change nothing, as the stock backend
  does.\footnote{\pgsrc{pingora-cache/src/memory.rs}{397}{\texttt{MemCache::purge}}
  returns \texttt{PurgeOutcome::NotFound} for an identified exact target,
  ``because it cannot safely match identified entries''.}

  For the purge path \texttt{pingora-proxy} itself drives, that key is
  always a primary: \texttt{HttpCache::purge} panics unless the request is
  still in its initial \texttt{CacheKey}
  phase,\footnote{\pgsrc{pingora-cache/src/lib.rs}{1813}{\texttt{HttpCache::purge}}
  matches on \texttt{self.phase} and panics for any phase other than
  \texttt{CacheKey} (\pgsrc{pingora-cache/src/lib.rs}{1844}{L1844}); only
  \pgsrc{pingora-cache/src/lib.rs}{1564}{\texttt{cache\_vary\_lookup}} ever
  sets a variance on the key, and it runs only from the \texttt{CacheKey} or
  \texttt{Stale} phase of an ordinary lookup --- a call sequence a purge
  request never enters
  (\pgsrc{pingora-proxy/src/proxy_purge.rs}{78}{\texttt{HttpProxy::proxy\_purge}}
  calls \texttt{purge()} directly, with no preceding lookup).} so the key it
  purges is \texttt{to\_compact()} of the key set for the request, before any
  lookup could narrow it to a variant slot, and its combined hash is
  \texttt{key.primary}
  itself\footnote{\pgsrc{pingora-cache/src/key.rs}{61}{\texttt{CacheHashKey::combined\_bin}}
  returns \texttt{primary\_bin()} unchanged when \texttt{variance\_bin()} is
  \texttt{None}.} (\reqref{FRE-COMPLIANCE.PURGE-01}). That is an observation
  about this one caller, not an assumption the backend relies on: the public
  \texttt{HttpCache::spawn\_async\_purge}\footnote{\pgsrc{pingora-cache/src/lib.rs}{1852}{\texttt{HttpCache::spawn\_async\_purge}}
  accepts any enabled phase and purges \texttt{key.to\_compact()}, which
  carries the variance once a \texttt{Vary} lookup has set one.} can present a
  varianced key, and resolving by the combined hash then purges exactly that
  variant rather than its primary.

  A purge addressed to a primary never touches a \texttt{Vary} variant,
  which lives at an unrelated shard path the primary key cannot derive
  (\reqref{FRE-COMPLIANCE.PURGE-01}), nor a per-writer temporary for the
  same key, which is a distinct directory entry identified by its own
  writer-identity and \texttt{.tmp} suffix (\reqref{FDS-STORAGE.PUBLISH-01}) that a
  purge addressed to the bare canonical name never names or opens.

  The purge operation itself is an \texttt{unlink} of that canonical path,
  performed while holding the key's publication lock
  (\reqref{FDS-STORAGE.METAUPDATE-01}). The backend answers with a
  \texttt{PurgeOutcome} rather than a boolean.\footnote{\pgsrc{pingora-cache/src/storage.rs}{128}{\texttt{PurgeOutcome}}.}
  When nothing exists at that path, and no store of the key is in flight
  for an invalidation to end (below), the backend shall report
  \texttt{NotFound}, not an error --- \texttt{unlink}'s own \texttt{ENOENT} is
  already exactly that signal, so
  idempotency and the absent-key case (\reqref{FRE-COMPLIANCE.PURGEABSENT-01}) need
  no separate bookkeeping. IF the \texttt{unlink} fails for any other
  reason, THEN the purge shall fail with an error, leaving the record in
  place. When the path existed, or an invalidation ended a store of the key
  in flight, the backend shall report \texttt{Purged}, carrying no entry
  identifier, since the backend assigns
  none.\footnote{Returning \texttt{Purged(None)} is what makes
  \pgsrc{pingora-cache/src/lib.rs}{1896}{\texttt{eviction.remove}} address the
  same key-only entry the backend admitted; an identifier the backend never
  issued would leave the entry tracked.}

  An invalidation shall be answered only once it is durable
  (\reqref{FRE-COMPLIANCE.PURGEDURABLE-01}): after releasing the lock, the
  backend shall synchronise the shard directory to stable storage
  (\reqref{ADR-PERSISTENCE.DURABILITY-01}), and only then report the
  outcome. It does so also when the path did not exist: a purge made to
  evict may have unlinked the record just before without synchronising, and
  a crash undoing that would otherwise bring the object back after the
  invalidation reported it absent. A shard directory that does not exist
  has nothing to synchronise. IF that synchronisation fails, THEN the purge
  shall fail with an error, since its answer might not survive a host
  crash. A purge made to evict (\texttt{PurgeType::Eviction}), as
  \texttt{HttpCache} purges the victims its eviction manager returns and as
  the restart scan purges its own (\reqref{FDS-EVICTION.STARTUP-01}),
  shall not synchronise: an eviction a host crash undoes brings back a valid
  object, which the restart scan admits, and evicts again if the cache is
  over capacity. The record the \texttt{unlink} removes is held open across
  it and closed last, after any synchronisation and outside the lock:
  closing it frees its storage, which can take a long time for a large
  record. Once the backend
  is read-only (\reqref{FDS-COMPAT.READONLY-01}), every purge shall fail with
  an error without touching the file (\reqref{FRE-COMPAT.HOTUPGRADE-01}).

  \textbf{Racing writers.} A same-key store in flight publishes by renaming
  its own per-writer temporary over the canonical path at completion, under
  the key's publication lock (\reqref{FDS-STORAGE.PUBLISH-01}). An
  invalidation shall end every store of the key in flight, under the same
  lock (\reqref{FRE-COMPLIANCE.PURGE-01}): it marks them ended in the
  in-flight registry, so that later lookups pass over them
  (\reqref{FDS-STORAGE.PARTIAL-01}), and each one's \texttt{finish},
  deciding under the lock, removes its temporary instead of renaming it and
  reports an error (\reqref{FDS-STORAGE.PUBLISH-01}). A store that published
  before the invalidation took the lock has its record removed by the
  \texttt{unlink}, so no store that began before an invalidation lands
  after it; a store that begins after it is not affected. When the key holds
  no record but a store in flight, the invalidation ends the store and
  reports \texttt{Purged(None)}, as the stock backend reports a purge that
  removed only a store in
  flight.\footnote{\pgsrc{pingora-cache/src/memory.rs}{405}{\texttt{MemCache::purge}}
  reports \texttt{Purged} when it removed a store in flight, a stored
  response, or both.} The purge never names or opens a writer's temporary:
  the writer removes it, as it does after any store that ends unpublished.
  The stock backend's purge instead removes the store's entry from its
  in-flight map, and the writer then panics on the
  \texttt{unwrap()}s\footnote{\pgsrc{pingora-cache/src/memory.rs}{267}{\texttt{MemMissHandler::finish}}
  unwraps its own lookup in \texttt{self.temp} at
  \pgsrc{pingora-cache/src/memory.rs}{273}{L273} and
  \pgsrc{pingora-cache/src/memory.rs}{275}{L275}; a concurrent
  \texttt{purge} racing ahead of it removes that same entry
  (\pgsrc{pingora-cache/src/memory.rs}{403}{L403}).} in its own
  \texttt{finish} --- tolerable only because that struct is explicitly
  documented as ``for testing only, not for production
  use.''\footnote{\pgsrc{pingora-cache/src/memory.rs}{71}{\texttt{MemCache}'s
  own doc comment}: ``For testing only, not for production use.''} Here the
  writer's \texttt{finish} completes, and reports the purge.

  A purge made to evict leaves stores in flight alone: a racing store's
  \texttt{rename} and the eviction's \texttt{unlink}, each made under the
  lock, land in either order, and whichever is last determines the outcome
  --- the last-writer-wins convergence \reqref{FRE-COMPLIANCE.REPLACE-01}
  requires among racing stores. A metadata update or expiry is no such
  racer: it opens the record under the same publication lock, so it acts on
  the record before the purge removes it, or finds the key absent after
  (\reqref{FDS-STORAGE.METAUPDATE-01}).

  \textbf{Safe under a racing reader.} A reader already streaming the
  canonical file holds it open from before the purge arrives
  (Figure~\ref{fig:sds-objectfile-lookup-seq}); POSIX \texttt{unlink}
  removes only the directory entry, and ``if the name was the last link to
  a file but any processes still have the file open, the file will remain
  in existence until the last file descriptor referring to it is
  closed.''\footnote{\href{https://man7.org/linux/man-pages/man2/unlink.2.html}{Linux
  manual page: \texttt{unlink(2)}}.} The reader therefore runs to the end
  of the body it already opened undisturbed, and the space it reads from is
  not reused for anything else while that descriptor remains open,
  satisfying \reqref{FRE-COMPLIANCE.READERISOLATION-01}
  without any reference-counting of the backend's own: the kernel already
  does it.

  Figure~\ref{fig:fds-purge-seq} drills into the \texttt{purge(target,
  Invalidation)} call Figure~\ref{fig:purge-seq} shows opaque to
  \texttt{Storage}. Variant records left behind once their primary is gone
  are outside this operation's own reach (\reqref{FRE-COMPLIANCE.PURGE-01});
  eviction removes them like any other record
  (\reqref{ADR-STORAGE.RECLAIM-01}).

  \begin{figure}[H]
    \centering
    \begin{adjustbox}{max width=0.9\textwidth}
    \begin{sequencediagram}
      \newinst[0.0]{HttpCache}{HttpCache}
      \newinst[2.9]{Storage}{Storage}
      \newinst[6.7]{FileSystem}{File System}
      \begin{call}{HttpCache}{purge(target, purge\_type)}{Storage}{PurgeOutcome}
        \begin{sdalt}{alt}{[Exact(Identified)]}
          \postlevel
          \begin{callself}{Storage}{{\small return NotFound --- no generation this backend issued}}{}
          \end{callself}
        \sdpart{[Active or Exact(KeyOnly)]}
          \postlevel
          \postlevel
          \begin{callself}{Storage}{{\small \shortstack{compute canonical path from\\hex(target.key().combined()) (\reqref{SDS-LAYOUT.SHARDING-01})}}}{}
          \end{callself}
          \begin{call}{Storage}{unlink(path) under the publication lock}{FileSystem}{ok / ENOENT}
          \end{call}
          \begin{sdalt}{opt}{[Invalidation]}
            \postlevel
            \begin{callself}{Storage}{{\small under the lock: mark the key's stores in flight ended}}{}
            \end{callself}
            \begin{call}{Storage}{fsync(shard directory)}{FileSystem}{}
            \end{call}
          \end{sdalt}
          \postlevel
          \begin{callself}{Storage}{{\small return Purged(None) if unlinked or a store was ended, else NotFound}}{}
          \end{callself}
        \end{sdalt}
      \end{call}
    \end{sequencediagram}
    \end{adjustbox}
    \caption{Purge, drilling into the call Figure~\ref{fig:purge-seq} shows
    opaque to \texttt{Storage}. A target naming a generation identifier is
    rejected unseen, since this backend issues none
    (\reqref{ADR-COMPAT.ENTRYID-01}); otherwise the canonical path is computed
    directly from the target key's combined hash, never scanned or looked up,
    and a single \texttt{unlink} either removes it or reports
    \texttt{ENOENT}, which the backend folds into \texttt{NotFound}. An
    invalidation also marks the key's stores in flight ended, under the lock,
    and then synchronises the shard directory before it reports; a purge made
    to evict does neither --- the \emph{opt} fragment's other branch. No
    per-writer temporary, and no file of any other key, is named or touched by
    this call.}\label{fig:fds-purge-seq}
  \end{figure}

  \reqdiaglink{overview}{\textit{View the domain overview.}}

  \reqexpect{verify}{TC-STORAGE.PURGE-01}{Verify that purge unlinks only the
  canonical path derived from the presented target key's combined hash (a
  primary or a variant), for an Active target and a key-only Exact target
  alike; that an identified Exact target is reported not-found with nothing
  unlinked; that purge is idempotent and reports not-found for an absent key;
  that an invalidation ends a same-key store in flight, which then
  publishes nothing while the request writing it and a reader already
  streaming it still read its body, and reports a key holding only such a
  store purged, while a purge made to evict leaves the store to land; that
  it never names a writer's temporary nor disturbs a concurrent reader's
  already-open descriptor; that an invalidation synchronises the
  directory before it reports, also when it found nothing, while a purge
  made to evict does not; that a removal is reported carrying no entry
  identifier; that a purge
  whose unlink fails reports an error and leaves the record in place; and
  that it fails once the backend is read-only.}
  \reqcover{satisfy}{IMP}

  \refine{ADR-STORAGE.OBJECTFILE-01}
  \refine{ADR-PERSISTENCE.DURABILITY-01}
  \satisfy{FRE-COMPLIANCE.PURGE-01}
  \satisfy{FRE-COMPLIANCE.PURGEABSENT-01}
  \satisfy{FRE-COMPLIANCE.PURGEDURABLE-01}
  \satisfy{FRE-COMPLIANCE.READERISOLATION-01}
  \trace{FDS-STORAGE.PUBLISH-01}
  \trace{FDS-STORAGE.PARTIAL-01}
  \trace{FRE-COMPLIANCE.REPLACE-01}
  \trace{SDS-LAYOUT.SHARDING-01}
  \trace{ADR-STORAGE.RECLAIM-01}
\end{requirement}

\begin{requirement}{FDS-STORAGE.EXPIRE-01}[draft][functionalRequirement]{Expire by updating the metadata}

  \localreqdiagram*{fdsexpire01}{\reqid}{\reqdesc}

  An expiry is keyed off a \texttt{PurgeTarget} exactly as a purge is, and the
  backend shall resolve it exactly as \reqref{FDS-STORAGE.PURGE-01} resolves
  one: an \texttt{Active} target, or an \texttt{Exact} target whose
  \texttt{CacheEntryKey} is \texttt{KeyOnly}, to the canonical path of that
  key; an \texttt{Exact} target carrying an \texttt{Identified} entry key to
  \texttt{NotFound}, since this backend issues no entry identifiers
  (\reqref{ADR-COMPAT.ENTRYID-01}). No other file is named or opened.

  \textbf{The operation.} Holding the key's publication lock
  (\reqref{FDS-STORAGE.METAUPDATE-01}), the backend shall open the record,
  read its stored metadata, derive the expired metadata from it by applying
  \texttt{CacheMeta::expire\_at} at the start of the current second, and
  write the result in place, as a metadata update that fits is written
  (\reqref{ADR-STORAGE.EXPIRE-01}), without reading or writing the body.
  The result always fits the record's region, whatever headroom it has
  left: \texttt{expire\_at} changes only the freshness deadline in the
  internal-metadata segment, and only ever lowers it, so the deadline's
  seconds take no more bytes than before and its nanoseconds, zero, take
  the fewest; and the backend writes the header segment back as stored,
  byte for byte. The expired metadata is therefore never longer than the
  stored metadata.

  Taking the lock for the whole read-derive-write sequence is what keeps two
  concurrent expiries, or an expiry and a revalidation's
  \texttt{update\_meta}, from clobbering one another: neither can read
  metadata the other is between reading and rewriting. This is why the backend
  derives the expired metadata from what is on disk rather than from anything
  the caller holds --- the caller of an expiry holds no \texttt{CacheMeta} at
  all.

  \textbf{Outcomes.} The backend shall report \texttt{PurgeOutcome::Expired}
  once the expired metadata is written, and so visible to lookups;
  \texttt{NotFound} when no record exists at the canonical path --- the
  absent-key case (\reqref{FRE-COMPLIANCE.PURGEABSENT-01}), which is not an
  error --- and an error when the record exists but cannot be read, decoded,
  or written. An expiry never reports \texttt{Purged}. \texttt{Expired}
  tells \texttt{HttpCache} the entry is still stored, so it stays in the
  eviction manager's accounting, which is correct: the record still occupies
  its full length on
  disk.\footnote{\pgsrc{pingora-cache/src/lib.rs}{1901}{\texttt{purge\_impl}}
  makes no eviction call for an \texttt{Expired} outcome.} A record
  that fails the read-time validation \reqref{SDS-LAYOUT.OBJECTFILE-01}
  specifies --- damaged, truncated, of another format version, or withdrawn
  --- is reported \texttt{NotFound}, matching what a lookup of that same
  record would report, and is left as it is for
  \reqref{FDS-STORAGE.RECLAIM-01}.

  \textbf{Races.} The record is opened under the lock, so a store or purge
  that lands before the expiry takes the lock is seen: the expiry acts on the
  record the key then holds, or reports \texttt{NotFound}. A same-key store
  already in flight is unaffected: it publishes by renaming over the canonical
  name when it completes, replacing the expired record with a fresh object,
  which is the behaviour \reqref{FRE-COMPAT.EXPIRE-01} states; an
  invalidating purge, unlike an expiry, ends such a store
  (\reqref{FDS-STORAGE.PURGE-01}). A reader already streaming the record
  keeps its open descriptor
  and is undisturbed, since the expiry does not touch the body
  (\reqref{FRE-COMPLIANCE.READERISOLATION-01}).

  Once the backend is read-only (\reqref{FDS-COMPAT.READONLY-01}), every
  expiry shall fail with an error without changing anything on disk
  (\reqref{FRE-COMPAT.HOTUPGRADE-01}).

  \begin{figure}[H]
    \centering
    \begin{adjustbox}{max width=0.95\textwidth}
    \begin{sequencediagram}
      \newinst[0.0]{HttpCache}{HttpCache}
      \newinst[1.0]{Storage}{Storage}
      \newinst[7.3]{FileSystem}{File System}
      \begin{call}{HttpCache}{expire(target)}{Storage}{PurgeOutcome}
        \begin{sdalt}{alt}{[Exact(Identified)]}
          \postlevel
          \begin{callself}{Storage}{{\small return NotFound (\reqref{ADR-COMPAT.ENTRYID-01})}}{}
          \end{callself}
        \sdpart{[Active or Exact(KeyOnly)]}
          \postlevel
          \begin{call}{Storage}{\shortstack{under the publication lock: open the canonical path,\\read header + metadata region}}{FileSystem}{CacheMeta / ENOENT}
          \end{call}
          \begin{sdalt}{opt}{[record valid]}
            \postlevel
            \begin{callself}{Storage}{{\small expire\_at(start of the second) --- internal segment only, body untouched}}{}
            \end{callself}
            \begin{call}{Storage}{write the segments in place, then the header}{FileSystem}{}
            \end{call}
            \postlevel
            \begin{callself}{Storage}{return Expired}{}
            \end{callself}
          \end{sdalt}
        \end{sdalt}
      \end{call}
    \end{sequencediagram}
    \end{adjustbox}
    \caption{Expiry, the second removal operation beside
    Figure~\ref{fig:fds-purge-seq}'s purge. The target is resolved the same
    way, but instead of unlinking the canonical name the backend derives
    expired metadata from what that record already stores and writes it back
    in place, as \reqref{FDS-STORAGE.METAUPDATE-01} writes an update that
    fits. The body is never read or copied, and the record survives, so the
    \texttt{304} that follows the revalidation serves the very bytes the cache
    already held. The read, the derivation, and the write are one critical
    section under the key's publication lock, so a concurrent expiry or
    metadata update cannot clobber it. The \emph{opt} fragment's other branch
    --- no record at the path, or one that fails validation --- reports
    \texttt{NotFound}.}\label{fig:fds-expire-seq}
  \end{figure}

  \reqdiaglink{overview}{\textit{View the domain overview.}}

  \reqexpect{verify}{TC-STORAGE.EXPIRE-01}{Verify that an expiry stops the
  object being served as a fresh hit while leaving its body byte-identical and
  readable; that the stored serve-stale windows are not closed by it; that it
  is idempotent; that it reports not-found for an absent key, for an
  identified Exact target, and for a record that fails the read-time
  validation, a withdrawn one included, which it leaves untouched; that it
  writes its metadata in place, with headroom configured and without, also
  for a deadline on a whole second; that a concurrent store still replaces
  the expired record; and that it fails once the backend is read-only.}
  \reqcover{satisfy}{IMP}

  \refine{ADR-STORAGE.EXPIRE-01}
  \satisfy{FRE-COMPAT.EXPIRE-01}
  \satisfy{FRE-COMPLIANCE.PURGEABSENT-01}
  \trace{FDS-STORAGE.METAUPDATE-01}
  \trace{FDS-STORAGE.PURGE-01}
  \trace{FRE-COMPLIANCE.READERISOLATION-01}
\end{requirement}

\begin{requirement}{ADR-STORAGE.EXECUTOR-01}[draft][designConstraint]{File I/O on the blocking pool}

  \localreqdiagram*{adrexecutor01}{\reqid}{\reqdesc}

  \textbf{Context:} Every operation this backend performs is an ordinary
  blocking syscall --- \texttt{open}, \texttt{pread}, \texttt{pwrite},
  \texttt{fdatasync}, \texttt{rename}, \texttt{unlink} --- while every
  \texttt{Storage} method is an \texttt{async fn} that Pingora polls on a
  service worker. Pingora gives each service \emph{one} worker thread by
  default,\footnote{\pgsrc{pingora-core/src/server/configuration/mod.rs}{74}{\texttt{ServerConf::threads}}
  defaults to 1 (\pgsrc{pingora-core/src/server/configuration/mod.rs}{236}{L236}).}
  and the deliverable cannot assume an integrator raised it. A syscall run
  inline therefore stalls not just the request that issued it but every other
  connection that worker is driving.

  The cost is not hypothetical, and not confined to slow devices. A cache
  whose whole purpose is to hold more than fits in RAM
  (\reqref{HLR-PERSISTENCE.CAPACITY-01}) serves reads that miss the page cache
  as a matter of course, and a store's \texttt{fdatasync}
  (\reqref{ADR-PERSISTENCE.DURABILITY-01}) costs 1--5~ms on a consumer SSD by
  this specification's own figures. Blocking a sole worker for that long
  reorders what the proxy does next: measured against
  the stock \texttt{MemCache} of \texttt{pingora-cache} on one worker, an inline
  backend lost an upstream \texttt{103 Early Hints} on 30 of 40 requests where
  \texttt{MemCache} lost none, because the proxy was late to poll its upstream
  task channel and the interim response was batched together with the final
  one. That is observable HTTP behaviour differing from the reference backend,
  which \reqref{HLR-COMPLIANCE.RFC-01} rules out.

  \textbf{Decision:} Run each storage operation's blocking work on Tokio's
  blocking pool --- \texttt{spawn\_blocking} --- rather than on the worker
  that polls the \texttt{async fn}. The handoff is per \emph{operation}, not
  per syscall: one for a lookup, one for a store's opening, one per body
  chunk, one for \texttt{finish}, one for a purge, a metadata update, or an
  expiry. The synchronous code that does the work is unchanged; only where it
  runs is.

  \textbf{Alternatives considered:}
  \begin{itemize}
    \item \emph{Leave the syscalls inline and document a minimum thread
      count}: no code to write, and an integrator who sets
      \texttt{ServerConf::threads} high enough is unaffected. Rejected: it
      makes correctness against \reqref{HLR-COMPLIANCE.RFC-01} a deployment
      setting, and the default --- one thread --- is the failing case, so the
      out-of-the-box configuration would be the broken one.
    \item \emph{\texttt{tokio::task::block\_in\_place}}: hands the worker's
      other tasks to a replacement thread without moving the work. Rejected:
      it panics on a current-thread runtime, so it would make the backend
      depend on which runtime flavour the integrator's Pingora build uses.
    \item \emph{\texttt{tokio::fs}}: the same blocking pool with an async
      surface. Rejected: it offloads one syscall at a time, so a store would
      pay several handoffs where one suffices, and its API does not cover the
      positional reads and writes (\texttt{pread}/\texttt{pwrite}) this
      layout depends on.
    \item \emph{\texttt{io\_uring}}: genuinely asynchronous file I/O, no pool
      and no handoff. Rejected for this version: a new dependency, a Linux
      version floor above what Pingora itself requires, and a second I/O path
      to keep correct --- machinery \reqref{HLR-SCOPE.V1-01} asks not to build
      when a pool that ships with the runtime already closes the gap.
  \end{itemize}
  Appendix~\ref{app:executor} records an investigation left to a later
  version: running the file I/O on the workers instead, \texttt{io\_uring},
  and the file systems that suit each.

  \textbf{Consequences:}
  \begin{itemize}
    \item[+] A storage operation never holds a service worker, so a slow
      device delays the request that touched it and nothing else.
    \item[+] Correct on the default Pingora configuration, with no requirement
      on the integrator.
    \item[--] One thread handoff per storage operation, which is charged
      against every lookup and every body chunk whether or not the data was
      in the page cache; on Linux, \reqref{ADR-STORAGE.PAGECACHE-01} keeps
      the reads the page cache answers on the worker.
    \item[--] The blocking pool is shared with everything else in the process
      that uses it, and its size is the runtime's to decide, not this
      backend's.
    \item[--] The restart scan is the one path this does not cover:
      \reqref{FDS-EVICTION.STARTUP-01}'s \texttt{recover()} borrows the
      eviction manager for the length of the call, so its work cannot be sent
      to another thread. It bounds the stall by yielding instead.
  \end{itemize}

  \reqdiaglink{overview}{\textit{View the domain overview.}}

  \reqexpect{refine}{FDS-STORAGE.EXECUTOR-01}{Specify which work moves to the
  blocking pool, what stays on the worker, and how the restart scan bounds
  the stall it cannot move.}

  \refine{ADR-PERSISTENCE.MECHANISM-01}
  \satisfy{HLR-COMPLIANCE.RFC-01}
  \trace{ADR-PERSISTENCE.DURABILITY-01}
  \trace{FRE-COMPAT.STREAMDELIVERY-01}
\end{requirement}

\begin{requirement}{FDS-STORAGE.EXECUTOR-01}[draft][functionalRequirement]{Storage calls never block the worker}

  \localreqdiagram*{fdsexecutor01}{\reqid}{\reqdesc}

  WHILE a \texttt{Storage} call, a \texttt{HandleHit} read, or a
  \texttt{HandleMiss} write is in progress, the backend shall leave the
  executor thread that polled it free to drive other tasks
  (\reqref{ADR-STORAGE.EXECUTOR-01}).

  \textbf{What moves.} Each of the following runs as one unit on the blocking
  pool, so an operation costs a single handoff rather than one per syscall:
  \begin{itemize}
    \item \texttt{lookup} --- opening the record, validating it, and
      deserialising its metadata (\reqref{SDS-LAYOUT.OBJECTFILE-01}), for a
      published record and for a partial hit's temporary alike
      (\reqref{FDS-STORAGE.PARTIAL-01});
    \item \texttt{get\_miss\_handler} --- creating the per-writer temporary
      and writing the record's fixed part (\reqref{FDS-STORAGE.PUBLISH-01});
    \item \texttt{write\_body} --- writing a chunk to the temporary;
    \item \texttt{finish} --- patching the header, the \texttt{fdatasync}
      unless stores are unsynchronised
      (\reqref{FDS-PERSISTENCE.CRASHDISCARD-01}), and the rename
      (\reqref{ADR-PERSISTENCE.DURABILITY-01}), then
      dropping the miss handler, whose close frees the store's file if the
      file has lost its name meanwhile;
    \item \texttt{read\_body} --- each positional read, on both the published
      and the partial-hit path;
    \item \texttt{purge} --- the unlink and the shard-directory
      \texttt{fsync} (\reqref{FDS-STORAGE.PURGE-01});
    \item \texttt{update\_meta} and \texttt{expire} --- reading the record and
      writing it back in place, or withdrawing or removing it
      (\reqref{FDS-STORAGE.METAUPDATE-01}, \reqref{FDS-STORAGE.EXPIRE-01}).
  \end{itemize}
  On Linux, \reqref{FDS-STORAGE.PAGECACHE-01} keeps on the worker the lookup
  of a published record, and the reads of its body, when no call has to
  wait for the device.

  \textbf{What stays.} Work that touches no file stays on the worker, because
  a handoff would cost more than it saves: hashing the cache key, serialising
  and deserialising \texttt{CacheMeta}, folding a chunk into the running
  CRC-32C, and every consultation of the in-memory structures --- the
  in-flight registry, the shard-directory bitset, and the publication lock's
  stripe table. A call therefore does its CPU work, hands the file work over,
  and resumes: the borrowed \texttt{CacheKey} and \texttt{CacheMeta} never
  have to cross the boundary, since only owned bytes do.

  A partial hit's wait for the next chunk also stays on the worker: it is a
  notification, not a syscall (\reqref{FDS-STORAGE.PARTIAL-01}), and blocking
  a pool thread on it would hold that thread for as long as the store takes.

  Two syscalls stay deliberately. The first is the \texttt{unlink} that
  removes an aborted store's temporary when its miss handler is dropped. A destructor
  cannot await, and the temporary's name must be gone when the drop returns
  (\reqref{FRE-COMPAT.ABORTSPACE-01}). The unlink is cheap: the handler still
  has the file open, so it removes only the name. The storage is freed when
  that handle is closed, which can take a long time for a large file, so
  the handler gives the handle to the pool to close; should the process
  exit first, the kernel closes it, and nothing is left behind.

  The second is an \texttt{fstat} when a hit handler is dropped holding the
  last descriptor of the file it read. Closing that descriptor frees the
  file if the file has lost its last name: a record replaced, purged or
  evicted while the hit read it, or a partial hit's temporary removed with
  its store. The \texttt{fstat} tells: a file that still has a name is
  closed in place, since its close frees nothing, and one that has none is
  given to the pool to close, without waiting. Handing every hit's close to
  the pool would cost the worker more than the close itself. The
  \texttt{fstat} reads the open file's inode, which is in memory, as the
  cache directory is on a local file system.

  \textbf{The publication lock.} The per-key lock and the read-only gate
  (\reqref{FDS-STORAGE.METAUPDATE-01}, \reqref{FDS-COMPAT.READONLY-01}) are
  taken and released entirely within one blocking unit, so neither is ever
  held across an \texttt{await} --- the same rule as before, now enforced by
  the shape of the code rather than only by convention.

  \textbf{The restart scan.} \reqref{FDS-EVICTION.STARTUP-01}'s
  \texttt{recover()} borrows the eviction manager it admits into for the
  length of the call, so its admissions cannot be sent to another thread.
  It shall instead yield to the executor between directory entries, and
  between the records its ghost-entry check visits, so that a scan of a large
  tree never holds the worker for longer than one record's header, key, and
  metadata region take to read. What the scan removes --- a
  record or temporary it reclaims, an eviction victim --- it shall remove on
  the blocking pool, since freeing a large file can take far longer than
  that read.

  \reqdiaglink{overview}{\textit{View the domain overview.}}

  \reqexpect{verify}{TC-STORAGE.EXECUTOR-01}{Verify that no storage operation
  holds the executor thread for its duration, by measuring what a task
  co-resident on a single-worker runtime observes while cache traffic runs.}
  \reqcover{satisfy}{IMP}

  \refine{ADR-STORAGE.EXECUTOR-01}
  \satisfy{HLR-COMPLIANCE.RFC-01}
  \trace{FDS-STORAGE.PUBLISH-01}
  \trace{FDS-STORAGE.PARTIAL-01}
  \trace{FDS-STORAGE.PURGE-01}
  \trace{FDS-STORAGE.METAUPDATE-01}
  \trace{FDS-EVICTION.STARTUP-01}
\end{requirement}

\begin{requirement}{ADR-STORAGE.PAGECACHE-01}[draft][designConstraint]{Page-cache reads on the worker}

  \localreqdiagram*{adrpagecache01}{\reqid}{\reqdesc}

  \textbf{Context:} \reqref{ADR-STORAGE.EXECUTOR-01} runs every storage
  operation on Tokio's blocking pool, and counts the cost among its
  consequences: one thread handoff per operation, whether or not the data
  was in the page cache. Measured on 22~October~2026 with the benchmark of
  \reqref{TC-PERFORMANCE.HITRATE-01} (four workers, 64 connections), that cost
  dominates a hit on an object the page cache holds. A 1~KiB hit took
  116~microseconds of processor time against 59 with \texttt{MemCache}, 42 of
  the difference on pool threads, for an \texttt{open} and a read the page
  cache answered; the pool grew to about 100 threads; and hits ran at 48, 53
  and 62\,\% of \texttt{MemCache} for 1~KiB, 64~KiB and 1~MiB bodies, short
  of \reqref{HLR-PERFORMANCE.HITRATE-01} for the first two.

  What \reqref{ADR-STORAGE.EXECUTOR-01} rules out is a worker waiting for the
  device. A call the page cache answers does not wait, and Linux can report
  in advance that a call would: \texttt{openat2} with
  \texttt{RESOLVE\_CACHED} fails with \texttt{EAGAIN} unless every component
  of the path is in the kernel's lookup
  cache,\footnote{\href{https://man7.org/linux/man-pages/man2/openat2.2.html}{Linux
  manual page: \texttt{openat2(2)}}, \texttt{RESOLVE\_CACHED} (since Linux
  5.12).} and \texttt{preadv2} with \texttt{RWF\_NOWAIT} fails with
  \texttt{EAGAIN} rather than wait for data the page cache does not
  hold.\footnote{\href{https://man7.org/linux/man-pages/man2/preadv2.2.html}{Linux
  manual page: \texttt{preadv2(2)}}, \texttt{RWF\_NOWAIT} (since Linux 4.14).}

  \textbf{Decision:} On Linux, perform the lookup of a published record, and
  each read of a published record's body, on the worker that polls it,
  through calls that fail rather than wait: open the record with
  \texttt{openat2} and \texttt{RESOLVE\_CACHED}, and read it with
  \texttt{preadv2} and \texttt{RWF\_NOWAIT}. When one of these calls reports
  that it would wait, or that the kernel lacks the call or the flag, run the
  whole operation on the blocking pool, as \reqref{ADR-STORAGE.EXECUTOR-01}
  decides; so too when a lookup meets a checksum mismatch, whose re-read takes
  the key's publication lock (\reqref{SDS-LAYOUT.OBJECTFILE-01}). Every other
  operation stays on the pool: partial hits, stores, purges, metadata
  updates, expiries and the restart scan. On other systems every operation
  stays on the pool. A configuration option sends every operation to the
  pool on Linux as well; by default it does not.

  Under \reqref{HLR-SCOPE.V1-01} this takes the exception: the mechanism it
  adds --- a second path for reads, on the worker --- raised the hit rates
  \reqref{TC-PERFORMANCE.HITRATE-01} measures by 43, 40 and 83\,\% for 1~KiB,
  64~KiB and 1~MiB bodies, to 69, 75 and 113\,\% of \texttt{MemCache}.

  \textbf{Alternatives considered:}
  \begin{itemize}
    \item \emph{Cap the blocking pool}
      (\texttt{ServerConf::max\_blocking\_threads}): no code at all.
      Rejected: four to eight threads raised hits by 14--34\,\% but cut
      stores by 82--90\,\%, since each store's \texttt{fdatasync} holds a
      pool thread for 16--25~ms on a consumer SSD and stores complete only as
      fast as the pool lets their syncs run together; and the pool's size is
      the integrator's setting, shared with the rest of the process.
    \item \emph{A small pool of the backend's own for reads}, away from the
      pool stores fill. Rejected: every read still costs a handoff, and four
      threads --- the best cap measured --- left 1~KiB hits at 55\,\% of
      \texttt{MemCache}, short of \reqref{HLR-PERFORMANCE.HITRATE-01}, for
      one more component to build.
    \item \emph{\texttt{io\_uring}}: no pool and no handoff for any
      operation. Rejected for the reasons \reqref{ADR-STORAGE.EXECUTOR-01}
      gives; it would also replace the pool for every operation, where this
      decision moves only the reads that need not wait.
    \item \emph{Map records into memory}, and check with \texttt{mincore}
      that their pages are resident before copying them. Rejected: the kernel
      can evict a page between the check and the copy, and the copy then
      waits for the device on the worker.
    \item \emph{Read the body's first chunk in the lookup's handoff}: halves
      the handoffs of a hit on a body of up to 64~KiB. Not adopted: such a hit
      still costs one handoff, and this decision removes both.
  \end{itemize}

  \textbf{Consequences:}
  \begin{itemize}
    \item[+] A hit on a record the page cache holds makes no handoff, and the
      pool stays at a few threads beside the workers.
    \item[+] The worker still never waits for the device, so
      \reqref{FDS-STORAGE.EXECUTOR-01} holds: a call that would wait says so
      instead. The RFC 9111 differential on a single worker
      (\reqref{TC-COMPLIANCE.RFC-01}) found no new deviation.
    \item[--] Two read paths to keep equivalent, the worker's and the pool's.
    \item[--] Linux only: other systems keep every read on the pool, and a
      Linux kernel older than 5.12 keeps lookups there.
    \item[--] A read the page cache cannot answer pays one failed call before
      it reaches the pool.
    \item[--] A direct dependency on \texttt{rustix}, which makes the two
      calls without \texttt{unsafe} code.
  \end{itemize}

  \reqdiaglink{overview}{\textit{View the domain overview.}}

  \reqexpect{refine}{FDS-STORAGE.PAGECACHE-01}{Specify which reads run on the
  worker, through which calls, when the operation falls back to the blocking
  pool, and the option that sends every read to the pool.}

  \refine{ADR-STORAGE.EXECUTOR-01}
  \satisfy{HLR-PERFORMANCE.HITRATE-01}
  \trace{HLR-SCOPE.V1-01}
  \trace{HLR-COMPLIANCE.RFC-01}
\end{requirement}

\begin{requirement}{FDS-STORAGE.PAGECACHE-01}[draft][functionalRequirement]{Reads from the page cache}

  \localreqdiagram*{fdspagecache01}{\reqid}{\reqdesc}

  WHERE page-cache reads are enabled on Linux, WHEN a lookup reaches a
  published record or a hit reads a published record's body, the backend
  shall perform that operation on the worker that polls it through file calls
  that fail rather than wait for the storage device; an operation one of whose
  calls fails so runs on the blocking pool from its start, as
  \reqref{FDS-STORAGE.EXECUTOR-01} specifies
  (\reqref{ADR-STORAGE.PAGECACHE-01}).

  \textbf{The calls.} The lookup opens the canonical path
  (\reqref{SDS-LAYOUT.SHARDING-01}) read-only with \texttt{openat2} and
  \texttt{RESOLVE\_CACHED}, and reads the record's header, key and metadata
  region (\reqref{SDS-LAYOUT.OBJECTFILE-01}) with \texttt{preadv2} and
  \texttt{RWF\_NOWAIT}; a body read reads its chunk the same way. A read that
  returns less than it asked for is continued from where it stopped. A call
  \emph{fails rather than waits} when it returns \texttt{EAGAIN}; a kernel
  that lacks \texttt{openat2}, \texttt{RESOLVE\_CACHED} or
  \texttt{RWF\_NOWAIT} returns \texttt{ENOSYS}, \texttt{EINVAL} or
  \texttt{EOPNOTSUPP}, which count the same. \texttt{fstat}, which reads the
  open file's inode, and \texttt{close} are unchanged.

  \textbf{What runs on the pool instead.} The lookup goes to the pool, and
  starts there again from the open, when the open or a read of the record's
  fixed part fails rather than waits; when the record's checksums do not
  match, since the re-read under the key's publication lock
  (\reqref{SDS-LAYOUT.OBJECTFILE-01}) waits for an in-place update to end; and
  when a call fails with any other error, which the pool's attempt then
  reports with its context. A path that the kernel's lookup cache knows not to
  exist is a miss, answered on the worker, and so is a record that fails
  validation for any reason other than a checksum mismatch. A body read goes
  to the pool, and reads its whole chunk again, when its read fails for any
  reason. Partial hits (\reqref{FDS-STORAGE.PARTIAL-01}), stores, purges,
  metadata updates, expiries and the restart scan always run on the pool.

  \textbf{The option.} \texttt{FsCacheConfig} carries a switch for page-cache
  reads, on by default. Off, every lookup and body read runs on the pool as
  \reqref{FDS-STORAGE.EXECUTOR-01} specifies. On systems other than Linux the
  switch has no effect: every operation runs on the pool.

  \reqdiaglink{overview}{\textit{View the domain overview.}}

  \reqcover{verify}{TC}

  \refine{ADR-STORAGE.PAGECACHE-01}
  \satisfy{HLR-PERFORMANCE.HITRATE-01}
  \trace{FDS-STORAGE.EXECUTOR-01}
  \trace{SDS-LAYOUT.OBJECTFILE-01}
\end{requirement}

\begin{requirement}{ADR-STORAGE.RECLAIM-01}[draft][designConstraint]{Reclaim via the restart-time scan}

  \localreqdiagram*{adr6}{\reqid}{\reqdesc}

  \textbf{Context:} One kind of file is left behind by otherwise-correct
  operation of this design and never revisited by any keyed
  \texttt{Storage} call again: a per-writer temporary a crash orphans before
  its rename (\reqref{FDS-STORAGE.PUBLISH-01}, \reqref{FRE-COMPAT.STARTRECLAIM-01}).
  It is not discoverable through
  \texttt{lookup}, \texttt{purge}, or \texttt{update\_meta}: every
  \texttt{Storage} call pingora-cache itself makes is keyed to one specific
  object, never an enumeration of what the tree actually holds. The one
  place this specification already walks the whole tree is Eviction
  Recovery's own reconstruction scan (\reqref{FDS-EVICTION.STARTUP-01}),
  run once at every process start and explicitly precedented by
  \reqref{ADR-COMPAT.EVICTION-01} for a structurally similar problem: no
  snapshot, no delta log, no distinction between a clean shutdown and a
  crash, because only a fresh scan can be trusted regardless of how the
  previous run ended.

  A purge of a primary leaves records behind too: its \texttt{Vary}
  variants, whose keys \reqref{FDS-STORAGE.PURGE-01} establishes the backend
  is neither handed nor able to derive at the moment the primary is purged
  (\reqref{FRE-COMPLIANCE.PURGE-01}). Unlike a temporary, such a variant is
  a published object, which the eviction manager counts against the size
  budget and evicts like any other (\reqref{FRE-COMPAT.EVICTION-01}).

  \textbf{Decision:} Reclaim crash-orphaned temporaries as part of that same
  restart-time scan, rather than by any separate mechanism --- together with
  two kinds of canonical record the scan meets anyway: one that fails
  validation, which lookups report as a miss but deliberately leave in place
  (\reqref{SDS-LAYOUT.OBJECTFILE-01}), and one a metadata update withdrew
  (\reqref{FDS-STORAGE.METAUPDATE-01}). While Eviction Recovery walks the tree
  to reconstruct the eviction manager, the same walk also classifies and
  removes them, at the cost of no second full-tree traversal and no new
  trigger. Leave \texttt{Vary} variant records to eviction, whatever their
  primary holds: the scan admits them like any other record.

  \textbf{Alternatives considered:}
  \begin{itemize}
    \item \emph{Reclaim, in the same scan, the variants no lookup can
      reach}: those whose primary is absent, fails validation, or carries a
      \texttt{provenance()} later than the variant's own \texttt{created()},
      which is the test \texttt{pingora-cache} itself
      applies.\footnote{\pgsrc{pingora-cache/src/lib.rs}{1564}{\texttt{HttpCache::cache\_vary\_lookup}}
      sets \texttt{valid\_after} to the primary's \texttt{meta.provenance()}
      (\pgsrc{pingora-cache/src/lib.rs}{1576}{L1576}); the following
      \pgsrc{pingora-cache/src/lib.rs}{1489}{\texttt{cache\_lookup}} on the
      variant key discards any result whose own \texttt{meta.created()} is
      earlier (\pgsrc{pingora-cache/src/lib.rs}{1518}{L1518}).} Rejected:
      that test settles reachability only when every variance comes from
      the \texttt{Vary} handling in \texttt{pingora-cache} itself, which consults the
      primary first. An integrator's \texttt{cache\_key\_callback} can set a
      variance
      itself;\footnote{\pgsrc{pingora-cache/src/key.rs}{125}{\texttt{CacheKey::set\_variance\_key}}
      is public.
      \pgsrc{pingora-proxy/src/proxy_cache.rs}{97}{\texttt{HttpProxy::proxy\_cache}}
      looks a key that already carries a variance up as a variant slot, and
      \pgsrc{pingora-cache/src/lib.rs}{1388}{\texttt{update\_variance}}
      keeps that variance when \texttt{cache\_vary\_filter} returns the same
      one.} its responses are then stored in variant slots with no primary
      ever written, and lookups serve them, so the scan would delete every
      such object at each start. A variant no lookup can reach costs only
      space the size budget already counts, and eviction removes it like any
      object no request reaches.
    \item \emph{Reclaim a purge's variants synchronously, inside
      \texttt{purge} itself}: rejected on structural grounds, not just
      simplicity --- \reqref{FDS-STORAGE.PURGE-01} shows that the proxy's
      purge path only ever gives \texttt{purge} the primary's
      \texttt{CompactCacheKey}, and \texttt{pingora-cache} hands the backend
      neither a variant's own key nor any index from primary to variant. At
      that moment, without an index of the backend's own (the next
      alternative), there is no key to purge a variant \emph{by}.
    \item \emph{Record each primary's variants, and reclaim them inside
      \texttt{purge}}: when a variant is stored, append its hash to a small
      list kept beside its primary, so that purging the primary can unlink
      the listed variants at once. Rejected: a new on-disk structure per
      varied primary, which must stay consistent while variants of one
      primary are stored and replaced concurrently, across processes, and
      through crashes --- and the restart-time scan must exist anyway, for
      crash-orphaned temporaries and for eviction reconstruction. The
      space a variant whose primary is gone holds is already counted against
      the size budget, so leaving it to eviction costs no correctness.
    \item \emph{A separate periodic background thread}, scanning the tree
      on its own timer independently of restart: rejected as machinery
      \reqref{HLR-SCOPE.V1-01}'s simplest-adequate-design preference does
      not ask for --- a timer, an interval to tune, and a cancellation path
      for shutdown or hot-upgrade drain --- when the restart-time scan is
      already mandatory for eviction-state correctness
      (\reqref{FRE-COMPAT.EVICTION-01}) even with reclaim entirely absent,
      and needs no new machinery to extend.
    \item \emph{Reclaim a crash-orphaned temporary opportunistically}, the
      moment some other component happens to be looking at its directory:
      rejected for lack of such a moment --- \reqref{FDS-LAYOUT.STARTUP-01}'s
      shard-directory cache only ever tracks which second-level directories
      \emph{exist}, learned by listing the first-level directories once at
      construction; it never lists a shard's files, and no ordinary store
      operation is ever already enumerating a shard's files.
  \end{itemize}

  \textbf{Consequences:}
  \begin{itemize}
    \item[+] No new subsystem: no timer, no scheduler, no second tree walk,
      no shutdown-drain interaction to design --- extends a scan the
      backend must run regardless.
    \item[+] One trigger, one code path, matching
      \reqref{ADR-COMPAT.EVICTION-01}'s own precedent of never
      distinguishing a clean restart from a crash restart.
    \item[+] Compatible with a hot upgrade by construction:
      \reqref{FRE-COMPAT.HOTUPGRADE-01} already anticipates background
      reclamation running on the incoming side while the outgoing process
      drains, under the same open-file safety
      argument \reqref{FDS-STORAGE.PURGE-01} gives for purge --- unlinking
      a file leaves an already-open reader on it undisturbed.
    \item[--] Variant records a purge of their primary left behind stay on
      disk, across restarts too. Their space is counted against the size
      budget, and capacity-driven eviction removes them like any unrequested
      object once the cache nears that budget; a cache that never nears it
      keeps them, as leftover files in \reqref{HLR-EFFICIENCY.SPACE-01}'s
      sense. This specification accepts that as a v1 cost, matching the
      trade-offs \reqref{ADR-EFFICIENCY.COMPACTION-01} and
      \reqref{ADR-COMPAT.EVICTION-01} already record.
    \item[--] A record withdrawn or damaged while a process is up is only
      reclaimed at that process's \emph{next} restart, unless a store of its
      key replaces it first; a deployment wanting a tighter bound must
      restart or hot-upgrade more often. This specification does not
      mandate a schedule for that. To make what accumulates visible, each
      scan logs the number and bytes of the files it reclaimed, by class,
      and also reports them as metrics when the metrics feature is enabled
      (\reqref{ADR-OBSERVABILITY.METRICS-01}).
  \end{itemize}

  \reqdiaglink{overview}{\textit{View the domain overview.}}

  \reqexpect{refine}{FDS-STORAGE.RECLAIM-01}{Specify the concrete detection
  and removal procedure for each class of file the scan reclaims, riding Eviction Recovery's
  existing enumeration, and the per-class count and bytes of reclaimed files
  each scan reports.}

  \refine{ADR-PERSISTENCE.MECHANISM-01}
  \satisfy{FRE-COMPAT.STARTRECLAIM-01}
  \trace{ADR-COMPAT.EVICTION-01}
  \trace{FDS-EVICTION.STARTUP-01}
  \trace{FRE-COMPAT.HOTUPGRADE-01}
\end{requirement}

\begin{requirement}{FDS-STORAGE.RECLAIM-01}[draft][functionalRequirement]{Reclaiming orphaned files}

  \localreqdiagram*{fdsreclaim01}{\reqid}{\reqdesc}

  Eviction Recovery's scan (\reqref{FDS-EVICTION.STARTUP-01}) already visits
  every file in the tree once, reading each canonical record's header, key,
  and internal metadata to \texttt{admit} it. This procedure
  (\reqref{ADR-STORAGE.RECLAIM-01}) rides that same pass: for each directory
  entry it classifies by name and, for a canonical record, by validation,
  then either admits it as before or reclaims it --- no second walk of the
  tree, and no entry visited twice. Three classes are reclaimed:
  crash-orphaned temporaries, damaged, foreign or misplaced canonical
  records, and withdrawn records.

  \textbf{Crash-orphaned temporaries.} An entry with the \texttt{.tmp}
  extension carries its writer's lease identity in its name
  (\reqref{FDS-STORAGE.PUBLISH-01}). The procedure shall treat the writer as
  possibly running --- and leave the temporary untouched --- exactly when
  the writer's lease exists and the procedure cannot take its lock, or
  cannot tell (\reqref{ADR-STORAGE.WRITERID-01}). The process's own writers
  satisfy this test, since serving may already have begun while the scan
  runs: a lock held through one open file description refuses another, even
  in the same process. So does a live hot-upgrade peer. Every other
  temporary belongs to a writer that has stopped, and the procedure shall
  unlink it. A temporary whose name has another form is left alone. The test
  requires every process sharing the cache directory to run on this host,
  with the directory on a local file system, a deployment constraint the
  deliverable documents.

  \textbf{Stopped writers' leases.} After the walk, the procedure shall
  remove every lease, and every lease left under its staging name, whose
  lock it can take, holding the lock across the removal; a lease whose lock
  is held belongs to a live writer, or to one being created, and is left.
  Since a lease gets its final name only once it is locked, a live writer's
  lease is never removed. A temporary whose lease is missing counts as its
  writer's having stopped, so the order of the two removals does not
  matter.

  \textbf{Damaged or foreign canonical records.} A canonical record that fails
  the read-time validation \reqref{SDS-LAYOUT.OBJECTFILE-01} specifies ---
  including one of another format version --- cannot be admitted, since its key or freshness
  deadline cannot be trusted or decoded. Lookups report such records as
  misses but leave them in place, so the procedure shall reclaim them. A
  record whose checksums mismatch on every read may instead be an in-place
  metadata update caught between its two writes: before treating it as
  damaged, the procedure reads it once more holding the key's publication
  lock, as a lookup does (\reqref{SDS-LAYOUT.OBJECTFILE-01}), and admits it
  if that read finds it whole.

  \textbf{Misplaced records.} A canonical record found anywhere but at the
  path its stored key names (\reqref{SDS-LAYOUT.SHARDING-01}) --- put there by
  hand, or by a defect --- is never served: a lookup of its key reads the
  other path, and a lookup of the key this path belongs to finds a key
  mismatch (\reqref{SDS-LAYOUT.OBJECTFILE-01}). The procedure shall reclaim
  it, counted as damaged.

  \textbf{Withdrawn records.} A canonical record that a metadata update
  withdrew (\reqref{FDS-STORAGE.METAUPDATE-01}) holds no object: lookups
  report it as a miss and leave it in place until a store of its key
  replaces it. The procedure shall reclaim it, under a class of its own, so
  that its count is not mistaken for damage.

  \textbf{Variants.} A canonical record whose own stored key carries a
  variance (\texttt{Some(variance)} in its \texttt{CompactCacheKey}) is a
  secondary slot (\reqref{SDS-LAYOUT.OBJECTFILE-01}). The procedure shall
  admit it like any other valid record, whatever its sibling primary holds
  (\reqref{ADR-STORAGE.RECLAIM-01}). It never opens that primary: a lookup
  can reach a variant for which no primary was ever stored, and a variant no
  lookup can reach is evicted like any other record.

  \textbf{Entries it cannot read.} Nothing is known of a directory the
  procedure cannot list, or of an entry it cannot read, so it can be neither
  admitted nor reclaimed. Unless the directory or entry has vanished
  meanwhile, the procedure shall skip it and name it in a warning in the
  log.

  \textbf{Reclaiming safely under live traffic.} This scan may run while the
  process serves (\reqref{FDS-EVICTION.STARTUP-01}), and while an outgoing
  hot-upgrade peer, its backend already read-only, completes its last requests
  (\reqref{FDS-COMPAT.READONLY-01}). A live store can publish a fresh record at a
  path the procedure has just classified as reclaimable, at any moment ---
  its upstream fetch may have begun long before. The procedure shall
  therefore never unlink a canonical record by path alone: it notes the
  file's identity (device and inode) when it opens the entry to classify it,
  and unlinks it only while holding the key's publication lock
  (\reqref{FDS-STORAGE.METAUPDATE-01}), after confirming that the path still
  refers to that same file; otherwise it skips the entry. Because the read-only
  peer publishes nothing, the lock --- which this process's stores and purges
  also take --- covers every writer. A temporary needs no such check: nothing
  is ever renamed onto a temporary's name. One race remains harmless by
  construction: a file the peer or a reader already holds open is unaffected
  by the unlink, for the same POSIX reason \reqref{FDS-STORAGE.PURGE-01}
  relies on.

  \textbf{Reporting.} When the scan completes, or stops because the backend
  is read-only or its server is shutting down
  (\reqref{FDS-EVICTION.STARTUP-01}), the procedure shall log one summary
  line, saying which, and giving the number of directories and
  entries it skipped and the number and bytes of the files it reclaimed in
  each
  class --- \texttt{temporary}, \texttt{damaged}, and \texttt{withdrawn} ---
  and, when the backend's \texttt{metrics} feature is enabled (it is off by
  default, \reqref{ADR-OBSERVABILITY.METRICS-01}), add the same figures to two counters
  in the \texttt{prometheus} default registry, each labelled by class:
  \begin{itemize}
    \item \texttt{pingora\_durable\_reclaimed\_files\_total}, the
      number of files;
    \item \texttt{pingora\_durable\_reclaimed\_bytes\_total}, their
      bytes.
  \end{itemize}
  The backend registers both when it is built, with every class's series
  at zero, so that each is served from the start rather than from the first
  reclaim of its class (as \reqref{FDS-OBSERVABILITY.HEADROOM-01} does for
  its counter).

  \begin{figure}[H]
    \centering
    \begin{adjustbox}{max width=\textwidth}
    \begin{sequencediagram}
      \newinst[0.0]{EvictionRecovery}{Eviction Recovery}
      \newinst[0.5]{Storage}{Storage}
      \newinst[6.7]{FileSystem}{File System}
      \newinst[1.2]{EvictionManager}{EvictionManager}
      \begin{call}{EvictionRecovery}{enumerate()}{Storage}{}
        \begin{sdalt}{loop}{[for each directory entry, in any order]}
          \begin{sdalt}{alt}{[classify by name and validation]}
            \postlevel
            \begin{call}{Storage}{\shortstack{.tmp: writer's lease held?\\no: unlink --- yes: skip}}{FileSystem}{}
            \end{call}
          \sdpart{[canonical, fails validation, misplaced, or withdrawn]}
            \postlevel
            \begin{call}{Storage}{\shortstack{under the publication lock: same file?\\yes: unlink --- no: skip}}{FileSystem}{}
            \end{call}
          \sdpart{[canonical, valid (primary or variant slot)]}
            \begin{call}{Storage}{admit(key, size, fresh\_until)}{EvictionManager}{}
            \end{call}
          \end{sdalt}
        \end{sdalt}
        \postlevel
        \begin{call}{Storage}{\shortstack{each lease in writers/: lock free?\\yes: unlink --- no: skip}}{FileSystem}{}
        \end{call}
      \end{call}
    \end{sequencediagram}
    \end{adjustbox}
    \caption{Reclaim, riding the same \texttt{enumerate()} pass
    Figure~\ref{fig:fds-startup-seq} shows as one opaque call: each entry is
    classified once; a crash-orphaned temporary or a damaged, foreign,
    misplaced, or withdrawn record is unlinked --- a canonical record only after
    confirming, under its publication lock, that the path still names the
    file classified --- and never admitted; every other record, variants
    included, is admitted exactly as \reqref{FDS-EVICTION.STARTUP-01} already
    specifies. After the walk, the leases of stopped writers are
    removed.}\label{fig:fds-reclaim-seq}
  \end{figure}

  \reqdiaglink{overview}{\textit{View the domain overview.}}

  \reqexpect{verify}{TC-STORAGE.RECLAIM-01}{Verify that a temporary whose
  writer is not running (its lease missing, or present but unlocked) is
  reclaimed, and one whose writer is running is left; that a stopped
  writer's lease is removed and a running writer's kept; that a damaged,
  foreign-format, misplaced, or withdrawn record is reclaimed; that a variant is
  admitted whatever its primary holds; that a record replaced by a live store
  between classification
  and unlink survives; that a directory or entry it cannot list or read is
  skipped, named in a warning, and counted; and that the log summary, and
  the reclaim counters
  when the metrics feature is enabled, report each class.}
  \reqcover{satisfy}{IMP}

  \refine{ADR-STORAGE.RECLAIM-01}
  \refine{ADR-STORAGE.WRITERID-01}
  \satisfy{FRE-COMPAT.STARTRECLAIM-01}
  \trace{FDS-EVICTION.STARTUP-01}
  \trace{FDS-STORAGE.PUBLISH-01}
  \trace{FDS-STORAGE.PURGE-01}
  \trace{SDS-LAYOUT.OBJECTFILE-01}
  \trace{SDS-LAYOUT.SHARDING-01}
  \trace{FRE-COMPAT.HOTUPGRADE-01}
\end{requirement}

\begin{requirement}{ADR-EFFICIENCY.COMPACTION-01}[draft][designConstraint]{Compaction deferred for v1}

  \localreqdiagram*{adr3}{\reqid}{\reqdesc}

  \textbf{Context:} Each object is stored as a self-contained record --- its
  metadata followed by its body (\reqref{ADR-STORAGE.OBJECTFILE-01}). Giving each
  record its own file still pays block rounding and the per-file overhead once
  per object: block-granular allocation may waste almost a full block on a
  small object's tail, a cost that does not shrink with the object's size and
  works directly against the stakeholder space-efficiency goal
  (\reqref{HLR-EFFICIENCY.SPACE-01}, \reqref{REQ-EFFICIENCY.FRAGMENTATION-01}).

  \textbf{Decision:} Do not implement object compaction for this version.
  Closing this waste requires real, multi-part machinery --- a background
  trigger and threshold policy, an aggregate file format and its own
  internal index, a lookup path that checks two locations instead of one,
  relocation on revalidation, and reclaiming space after a purge --- against
  which \reqref{HLR-SCOPE.V1-01} prefers the simplest adequate design, even
  at the cost of peak space efficiency. This decision therefore conflicts
  with \reqref{HLR-EFFICIENCY.SPACE-01} and with the derived
  \reqref{REQ-EFFICIENCY.FRAGMENTATION-01}, which no other decision
  satisfies for this version: the waste described above is accepted, not
  eliminated, for as long as it stands. The resolution, for both conflicts,
  is to revisit compaction in a later version, once
  \reqref{HLR-SCOPE.V1-01}'s scope constraint no longer applies.

  \textbf{Alternatives considered:}
  \begin{itemize}
    \item \emph{Implement compaction now}: a deferred background step that
      moves a completed object's record, back-to-back with others, into a
      shared \emph{aggregate} file once its wasted space passes a
      configurable threshold. It would still require choosing the aggregate
      backend --- a custom format reusing
      \reqref{ADR-STORAGE.OBJECTFILE-01}'s record layout, or an embedded
      engine such as libmdbx or SQLite. Rejected: real,
      multi-part complexity --- the aggregate format and index, the
      two-location lookup, relocation and reclaim handling, and the
      backend choice itself --- for a version \reqref{HLR-SCOPE.V1-01}
      asks to keep simple, even though it would fully close
      \reqref{REQ-EFFICIENCY.FRAGMENTATION-01}.
    \item \emph{Do not implement compaction for this version} (chosen):
      every object keeps its own file; the block-rounding and per-file
      waste on small objects is accepted rather than eliminated.
  \end{itemize}

  \textbf{Consequences:}
  \begin{itemize}
    \item[+] No new subsystem: no aggregate format, no index, no
      background scheduler, no second lookup location, no purge-hole
      reclamation --- matches \reqref{HLR-SCOPE.V1-01}'s preference for the
      simplest adequate design.
    \item[+] Nothing beyond what \reqref{ADR-STORAGE.OBJECTFILE-01} already
      specifies needs building or trusting.
    \item[--] Block-granular allocation and per-file overhead keep scaling
      with the number of small cached objects, unresolved for this
      version.
    \item[--] \reqref{REQ-EFFICIENCY.FRAGMENTATION-01}'s obligation is not
      satisfied by this decision.
  \end{itemize}

  \reqdiaglink{overview}{\textit{View the domain overview.}}

  \reqterminal[Deferred: compaction is not implemented for this version, so
  nothing further refines, satisfies, or verifies this decision; revisit if
  HLR-SCOPE.V1-01's scope changes.]

  \refine{ADR-PERSISTENCE.MECHANISM-01}
  \conflictsWith{HLR-EFFICIENCY.SPACE-01}
  \conflictsWith{REQ-EFFICIENCY.FRAGMENTATION-01}
  \trace{ADR-STORAGE.OBJECTFILE-01}
\end{requirement}

% ---------------------------------------------------------------------
% Pingora compatibility: the eviction-reconstruction decision and the
% detailed design record that fixes its mechanism.
% ---------------------------------------------------------------------
\begin{requirement}{ADR-COMPAT.EVICTION-01}[draft][designConstraint]{Eviction-state reconstruction strategy}

  \localreqdiagram*{adr5}{\reqid}{\reqdesc}

  \textbf{Context:} \reqref{FRE-COMPAT.EVICTION-01} requires the
  \texttt{pingora-cache} eviction manager's view to stay consistent with the
  objects persisted across a restart, but the manager is driven by
  \texttt{HttpCache}, not by the \texttt{Storage} backend, so the backend
  cannot keep it synchronised on its own --- reconstructing it from the
  persisted set is one way to close that gap. Doing so can trade startup
  latency against the length of any inconsistency window, reliance on a
  custom \texttt{EvictionManager} versus the stock one, and how far the
  result departs from a pure drop-in \texttt{Storage} backend
  (\reqref{HLR-COMPAT.PINGORA-01}); the alternatives below weigh those
  trade-offs.

  \textbf{Decision:} Adopt \emph{file-system reconstruction as the sole
  mechanism, on every restart}: enumerate the persisted entries
  (\reqref{ADR-PERSISTENCE.MECHANISM-01}) and
  \texttt{admit}\footnote{\pgsrc{pingora-cache/src/eviction/mod.rs}{252}{\texttt{EvictionManager::admit}}.}
  each one into the manager, purging any victim an \texttt{admit} returns.
  No snapshot is taken, no delta log is maintained, and no distinction is
  made between a clean shutdown and a crash: reconstruction always follows
  the same procedure, and it admits every object with its record size, the
  same size a completed store reports (\reqref{FRE-COMPAT.EVICTION-01}).

  Whether the integrator begins serving before reconstruction completes
  (accepting inaccurate capacity accounting until it does) or waits for it
  (accepting startup latency proportional to the number of stored objects)
  is an integration choice this decision leaves open. The window of
  inaccuracy therefore ends exactly when the procedure completes; no fixed
  time bound is claimed for it. Running alongside live traffic creates three
  races, each closed by a rule the procedure follows:
  \begin{itemize}
    \item \emph{A stale size.} Live traffic may replace an object between
      the scan reading its record and admitting it. The scan therefore
      re-checks the record immediately before admitting, and uses the size
      it finds then --- or skips the entry if the object is gone, since the
      live store already admitted its replacement.
    \item \emph{A ghost entry.} A live eviction drops its victim from the
      manager at once, and removes the victim's record from storage only
      afterwards, in a task of its
      own,\footnote{\pgsrc{pingora-cache/src/lib.rs}{1073}{\texttt{HttpCache::finish\_miss\_handler}}
      admits the completed store (\pgsrc{pingora-cache/src/lib.rs}{1118}{L1118}),
      then spawns a task that purges each victim the manager returned
      (\pgsrc{pingora-cache/src/lib.rs}{1128}{L1128}).} so the scan may
      admit the victim again in between, and its entry would outlive the
      record. After enumeration, the procedure re-checks every object it
      admitted and removes the manager's entry for any that no longer
      exists. A purge leaves no such entry: \texttt{HttpCache} removes a
      purged object's entry itself, after the
      purge.\footnote{\pgsrc{pingora-cache/src/lib.rs}{1896}{L1896}, once
      the backend has reported the object purged.}
    \item \emph{An early victim.} An admission during the scan may return,
      as a victim, an object live traffic has just read. It is purged like
      any other victim: capacity accounting stays exact, and the object is
      simply fetched again when next requested.
  \end{itemize}

  \textbf{Alternatives considered:}
  \begin{itemize}
    \item \emph{Snapshot on clean shutdown, reconstruction
      only on crash}: on a clean shutdown, snapshot the manager with
      \texttt{save};\footnote{\pgsrc{pingora-cache/src/eviction/mod.rs}{303}{\texttt{EvictionManager::save}}.}
      at startup, if the snapshot file is present, \texttt{load}\footnote{\pgsrc{pingora-cache/src/eviction/mod.rs}{306}{\texttt{EvictionManager::load}}.}
      it and delete it durably before serving resumes, so its presence
      means exactly ``a durable, unconsumed snapshot exists'' and its
      absence means ``reconstruct.'' That still leaves real machinery: a
      shutdown-side lifecycle hook in Eviction Recovery
      (\reqref{FDS-EVICTION.STARTUP-01}) --- today there is only one, at
      process start --- a rule that a draining process never saves (its
      snapshot would miss what the incoming process stores meanwhile), tied
      to the read-only call \reqref{ADR-COMPAT.READONLY-01} introduces, and a
      startup branch choosing between loading and scanning.
      \reqref{HLR-SCOPE.V1-01} prefers the simplest adequate design even at
      the cost of peak space or I/O efficiency; a recurring per-restart scan
      is exactly such an efficiency cost, and the crash path already has to
      build the scan/admit/reconcile machinery regardless
      (\reqref{FDS-EVICTION.STARTUP-01}). Rejected: a second lifecycle hook,
      a save rule coupled to the read-only switch, and a second startup path are more
      machinery than the restart-time saving justifies.
    \item \emph{Save/load with log replay}: snapshot the manager and keep an
      append-only delta log between snapshots; replay the log at startup
      instead of scanning the file system. Rejected: adds delta-log maintenance
      and save-time reset machinery to the steady-state hot path, which the
      adopted approach avoids entirely.
    \item \emph{Lazy admission}: no scan at all; admit an object the first
      time a lookup finds it. Rejected: an object nobody requests again is
      never admitted, so it is never counted against the size budget and
      never evicted --- exactly the gap \reqref{FRE-COMPAT.EVICTION-01}
      exists to close.
  \end{itemize}

  \textbf{Consequences:}
  \begin{itemize}
    \item[+] Simplest possible mechanism: no snapshot, no delta log, no
      validity marker, no startup branch --- matches \reqref{HLR-SCOPE.V1-01}'s
      preference for the simplest adequate design.
    \item[+] One behaviour regardless of how the previous run ended: nothing
      to misclassify a restart as clean when it was not, or vice versa.
    \item[+] Reuses the stock \texttt{EvictionManager} only through
      \texttt{admit}; no custom manager, and no \texttt{save}/\texttt{load}
      pairing to keep synchronised.
    \item[--] Every restart pays a full file-system scan; starting to serve
      before it completes leaves capacity accounting inaccurate until it
      does, and waiting for it to finish instead trades that window for
      startup latency proportional to the number of stored objects ---
      unlike the snapshot alternative's clean-restart case, which could
      avoid both costs at once.
    \item[--] Admission order follows the scan, not past use, so the
      recency order the manager chose victims by before the restart is lost:
      right after a restart, a frequently used object can be evicted ahead of
      one nobody requests (accepted for this version,
      \reqref{FRE-COMPAT.EVICTION-01}).
    \item[--] Reconstruction is carried out by a separate component, Eviction
      Recovery, additional to the \texttt{Storage} trait implementation
      (\reqref{FDS-EVICTION.STARTUP-01}), so the deliverable is more than a pure
      drop-in \texttt{Storage} backend (\reqref{HLR-COMPAT.PINGORA-01}).
  \end{itemize}

  \reqdiaglink{overview}{\textit{View the domain overview.}}

  \reqexpect{refine}{FDS-EVICTION.STARTUP-01}{Specify the concrete
  reconstruction procedure, and the API surface that orchestrates it,
  realising this strategy.}

  \satisfy{FRE-COMPAT.EVICTION-01}
  \refine{ADR-PERSISTENCE.MECHANISM-01}
\end{requirement}

\begin{requirement}{ADR-COMPAT.ENTRYID-01}[draft][designConstraint]{No storage-assigned entry IDs}

  \localreqdiagram*{adrentryid01}{\reqid}{\reqdesc}

  \textbf{Context:} Pingora 0.9.0 lets a backend give the eviction manager an
  identity finer than the cache key. \texttt{HandleHit} and \texttt{HandleMiss}
  may each return a \texttt{CacheEntryId},\footnote{\pgsrc{pingora-cache/src/storage.rs}{331}{\texttt{HandleHit::entry\_id}}
  and \pgsrc{pingora-cache/src/storage.rs}{393}{\texttt{HandleMiss::entry\_id}},
  both defaulting to \texttt{None}.} which \texttt{HttpCache} pairs with the
  \texttt{CompactCacheKey} into the \texttt{CacheEntryKey} it admits, accesses,
  and removes. The facility exists for a backend that keeps several
  \emph{generations} of one key at once, so that eviction accounts for the exact
  stored object rather than for whichever one currently answers the key. The
  choice is all-or-nothing and must be consistent across both handles.

  This backend has no such generations to distinguish. One key names one
  canonical path (\reqref{SDS-LAYOUT.SHARDING-01}) and a store replaces whatever
  is there by renaming over it (\reqref{FDS-STORAGE.PUBLISH-01}); a per-writer
  temporary is not a readable generation, and the identifier is required to
  survive \texttt{finish}, so a temporary's identity could not serve as one.

  \textbf{Decision:} Return \texttt{None} from both \texttt{entry\_id} methods
  and keep eviction keyed on the \texttt{CompactCacheKey} alone. Every entry
  this backend admits is therefore a \texttt{CacheEntryKey::KeyOnly}, and every
  purge outcome it reports is \texttt{Purged(None)}
  (\reqref{FDS-STORAGE.PURGE-01}). An \texttt{Exact} purge target carrying an
  \texttt{Identified} entry key names a generation this backend never issued, so
  it is reported as not found rather than resolved to the key's canonical path.

  \textbf{Alternatives considered:}
  \begin{itemize}
    \item \emph{Assign an identifier per published record} (its inode number, or
      a counter recorded in the record header): rejected --- it buys accuracy
      only for a backend that retains generations, and this one does not. Here
      it would make accounting worse: a store that replaces a key would be
      admitted as a second entry, while the manager, which storage never tells
      of the replacement, would keep the first until it aged out, counting
      space the rename already freed. It would also cost an identity that must
      be read back on every hit, and survive a restart scan.
    \item \emph{Retain generations and identify them}, publishing each store
      under its own name and letting eviction remove them individually:
      rejected --- it is a different storage design, not an addition to this
      one, and \reqref{FRE-COMPLIANCE.REPLACE-01}'s last-writer-wins
      convergence over a single canonical name is what several other decisions
      here already rest on.
  \end{itemize}

  \textbf{Consequences:}
  \begin{itemize}
    \item[+] The restart scan and live traffic address the same entry without
      having to agree on an identifier the scan would have to recover from disk
      (\reqref{FDS-EVICTION.STARTUP-01}).
    \item[+] The backend's eviction accounting is exactly what it was under
      Pingora 0.8.1, where the manager took the cache key directly.
    \item[--] An eviction victim arrives as an \texttt{Exact} target and a purge
      request as an \texttt{Active} one, yet both resolve identically here, so
      the distinction the framework draws is one this backend can only
      degenerate.
    \item[--] Adopting generations later is a change to this decision and to the
      record layout, not an incremental addition.
    \item[--] An eviction purge can remove a newer object stored under the
      victim's key. \texttt{HttpCache} purges victims in a spawned task after
      \texttt{admit} chooses them;\footnote{\pgsrc{pingora-cache/src/lib.rs}{1128}{The
      miss handler's \texttt{finish} path} spawns the purges of the victims
      \texttt{admit} returned.} a store of that key completing in between is
      what the purge removes. The eviction manager then tracks an entry with no
      record behind it until it is evicted again, and the next lookup misses.
      It is rare, since a victim is the least recently used key, and costs one
      refetch; the stock \texttt{MemCache}, which assigns no identifiers
      either, behaves the same.
  \end{itemize}

  \reqdiaglink{overview}{\textit{View the domain overview.}}

  \reqcover{satisfy}{IMP}

  \satisfy{FRE-COMPAT.EVICTION-01}
  \refine{ADR-PERSISTENCE.MECHANISM-01}
  \trace{ADR-COMPAT.EVICTION-01}
  \trace{FDS-STORAGE.PURGE-01}
  \trace{FDS-EVICTION.STARTUP-01}
\end{requirement}

\begin{requirement}{ADR-COMPAT.READONLY-01}[draft][designConstraint]{Read-only on the upgrade signal}

  \localreqdiagram*{adrreadonly01}{\reqid}{\reqdesc}

  \textbf{Context:} Once the outgoing process receives the Pingora upgrade
  signal, its backend must stop changing the shared cache --- no
  publication, and purges and metadata updates refused --- while still
  serving reads (\reqref{FRE-COMPAT.HOTUPGRADE-01}); meanwhile the incoming
  process must reach a consistent eviction view
  (\reqref{FRE-COMPAT.EVICTION-01}). The \texttt{Storage} trait receives no
  such signal. Pingora does announce it to the process's own code: the
  server broadcasts \texttt{ExecutionPhase::GracefulUpgradeTransferringFds}
  when \texttt{SIGQUIT} arrives, before it sends the listening sockets to the
  new process, and both processes then accept connections for
  \texttt{CLOSE\_TIMEOUT} (5~s) before the old one stops listening and drains
  its sessions for up to \texttt{EXIT\_TIMEOUT} (300~s).\footnote{\pgsrc{pingora-core/src/server/mod.rs}{235}{\texttt{Server::watch\_execution\_phase}};
  the phase is sent at \pgsrc{pingora-core/src/server/mod.rs}{279}{L279},
  before the sockets are sent (\pgsrc{pingora-core/src/server/mod.rs}{290}{L290});
  \pgsrc{pingora-core/src/server/mod.rs}{66}{\texttt{CLOSE\_TIMEOUT}} and
  \pgsrc{pingora-core/src/server/mod.rs}{63}{\texttt{EXIT\_TIMEOUT}}.}

  \textbf{Decision:} The backend exposes an operation that makes it
  \emph{read-only}, which the integrator invokes when its process's server
  broadcasts \texttt{GracefulUpgradeTransferringFds}. The two words are kept
  apart: the \emph{process} drains --- it finishes the sessions it has
  accepted --- while its \emph{backend} is read-only, reading the shared
  cache but no longer changing it. From that call on, the backend
  publishes nothing (a store that completes has its temporary removed
  instead of renamed), fails purges and metadata updates with an error, and
  keeps serving lookups and reads. The incoming process starts its restart
  procedure only once it holds the listening sockets \emph{and} the
  upgrade's close timeout has elapsed: the outgoing process's read-only call ---
  delivered asynchronously, shortly after its phase broadcast --- has taken
  effect by then, so every object it published is already on disk before the
  incoming scan begins.

  \textbf{Alternatives considered:}
  \begin{itemize}
    \item \emph{No read-only switch: two independent writers for the whole
      window.}
      Rejected: objects the outgoing process publishes after the incoming
      scan are unknown to the incoming eviction manager until the next
      restart, violating \reqref{FRE-COMPAT.EVICTION-01}, and the two
      managers together can let disk use exceed the budget.
    \item \emph{An on-disk generation marker}: each backend records a
      generation at start and checks, before every publication or purge, that
      no newer one exists --- no integrator hook needed. Rejected: a new
      on-disk structure (\reqref{HLR-SCOPE.V1-01}), a file-system check on
      every write, and a window between the check and the rename that the
      marker cannot close.
    \item \emph{The incoming process defers all writes until the outgoing one
      exits.} Rejected: it runs uncached for up to \texttt{EXIT\_TIMEOUT} ---
      five minutes of full origin load after every upgrade.
  \end{itemize}

  \textbf{Consequences:}
  \begin{itemize}
    \item[+] The incoming process is the sole writer for its whole life after
      the close timeout, so the single-instance guarantees apply unchanged.
    \item[+] The integration is one call from a Pingora hook the server
      already provides.
    \item[--] An integrator that does not wire the call gets the two-writer
      behaviour rejected above; the deliverable documents the call as
      required for hot upgrades.
    \item[--] Purges sent to the outgoing process once its backend is
      read-only fail and must be retried; responses it stores from then on
      are served but not cached.
    \item[--] The incoming process's restart procedure starts at least 5~s
      after it takes over, lengthening the window of inaccurate accounting
      \reqref{ADR-COMPAT.EVICTION-01} already accepts.
  \end{itemize}

  \reqdiaglink{overview}{\textit{View the domain overview.}}

  \reqexpect{refine}{FDS-COMPAT.READONLY-01}{Specify the read-only operation
  and its integration with the execution-phase broadcast, the behaviour of each
  Storage call after it, and the incoming process's start condition for its
  restart procedure.}

  \refine{ADR-PERSISTENCE.MECHANISM-01}
  \satisfy{FRE-COMPAT.HOTUPGRADE-01}
  \trace{ADR-COMPAT.EVICTION-01}
\end{requirement}

% ---------------------------------------------------------------------
% Standards compliance: detailed, verifiable storage requirements
% refining HLR-COMPLIANCE.RFC-01.
%
% Scope: this component is a pingora-cache Storage backend, not a full
% HTTP cache. Decisions such as whether a response is cacheable, how
% freshness/age are computed, when to revalidate, and how Cache-Control
% request/response directives are interpreted live in pingora-cache.
% The requirements below therefore capture only what the *storage layer*
% must do so those RFC 9111 / RFC 9110 semantics survive a store -> retrieve
% cycle intact: faithful round-trip of metadata and body, support for the
% operations pingora-cache invokes (metadata update on validation, variant
% selection, purge), and refusal to serve partially-stored responses.
%
% Each requirement is written to be verifiable so it can be checked by the
% test-case specifications in the following stage.
% ---------------------------------------------------------------------
\begin{requirement}{FRE-COMPLIANCE.HEADERS-01}[draft][functionalRequirement]{Faithful header round-trip}

  \localreqdiagram*{rfr01}{\reqid}{\reqdesc}

  WHEN \texttt{pingora-cache} looks up a stored object, including after a
  process restart, the storage backend shall return the complete response
  head it was handed for that object --- the status code, the protocol
  version, and every response header field with its name and value
  preserved byte for byte and same-named fields in their relative order.

  ``Every header field'' covers any head \texttt{pingora-cache} hands over:
  fields \texttt{pingora-cache} does not itself recognise, opaque non-ASCII
  octets in a value, repeated field names (including
  \texttt{Set-Cookie}-style fields that must not be list-combined), empty
  field values, and a header set large enough to exceed any fixed-size
  metadata region the on-disk format reserves. Returning the head unchanged
  keeps the status, freshness, validation, and \texttt{Cache-Control} inputs
  \texttt{pingora-cache} relies on from being dropped or altered by storage.

  \textit{Basis: \rfcsec{9111}{3.1}[storing header fields];
  \rfcsec{9110}{5.2}[field lines], \rfcsec{9110}{5.3}[field order],
  \rfcsec{9110}{6.2}[control data], \rfcsec{9110}{6.3}[header fields].}

  \reqdiaglink{overview}{\textit{View the domain overview.}}

  \reqcover{verify}{TC}
  \reqcover{satisfy}{SDS}

  \deriveReqt{HLR-COMPLIANCE.RFC-01}
\end{requirement}

\begin{requirement}{FRE-COMPLIANCE.TRAILERS-01}[draft][functionalRequirement]{Trailer fields kept distinct}

  \localreqdiagram*{rfr02}{\reqid}{\reqdesc}

  The storage backend shall return no field it received as a trailer field
  to \texttt{pingora-cache} as a header field.

  The backend may persist trailer fields separately from the header fields
  or discard them; either meets the requirement. In the targeted
  \texttt{pingora-cache} version the \texttt{Storage} interface offers no
  trailer channel: the metadata handed to the backend carries header fields
  only, and the miss handler carries body bytes only, so the framework can
  never present a trailer to the backend and the interface itself discharges
  this obligation. The requirement still constrains the backend's own on-disk
  format --- which offers no way to fold trailer-originated fields into the
  stored header set --- and any future framework version that does carry
  trailers.

  Because no trailer can reach the backend through the targeted interface, no
  test can exercise this obligation; it is verified by \emph{inspection} of the
  \texttt{Storage} interface and of the on-disk record format, recorded in
  \reqref{TC-COMPLIANCE.TRAILERS-01}.

  \textit{Basis: \rfcsec{9111}{3.1}[``Caches MUST NOT combine trailer fields with
  header fields'']; \rfcsec{9110}{6.5}[trailer fields].}

  \reqdiaglink{overview}{\textit{View the domain overview.}}

  \reqcover{verify}{TC}
  \reqcover{satisfy}{SDS}

  \deriveReqt{HLR-COMPLIANCE.RFC-01}
\end{requirement}

\begin{requirement}{FRE-COMPLIANCE.BODY-01}[draft][functionalRequirement]{Faithful round-trip of the body}

  \localreqdiagram*{rfr03}{\reqid}{\reqdesc}

  The storage backend shall deliver, on every read of a stored object's body
  --- from its start, or from a position sought through the hit
  handler\footnote{\pgsrc{pingora-cache/src/storage.rs}{270}{\texttt{HitHandler::can\_seek}}
  and \pgsrc{pingora-cache/src/storage.rs}{284}{\texttt{HitHandler::seek}}.}
  --- exactly the corresponding bytes of the body written for that object, in
  order, with no truncation, padding, re-encoding, or other transformation.

  The requirement applies to bodies of every size: empty bodies, bodies at the
  backend's chunk- and block-size boundaries, and bodies spanning many chunks.
  Exact reads from a sought position matter because \texttt{pingora-cache}
  builds range responses from a sought reader, so those responses then carry
  exactly the stored bytes for the requested range.

  \textit{Basis: \rfcsec{9110}{8.1}[representation data];
  \rfcsec{9111}{3}[storing responses]; \rfcsec{9110}{14.2}[Range].}

  \reqdiaglink{overview}{\textit{View the domain overview.}}

  \reqcover{verify}{TC}
  \reqcover{satisfy}{SDS}

  \deriveReqt{HLR-COMPLIANCE.RFC-01}
\end{requirement}

\begin{requirement}{FRE-COMPLIANCE.FRESHNESS-01}[draft][functionalRequirement]{Preserve freshness and age inputs}

  \localreqdiagram*{rfr04}{\reqid}{\reqdesc}

  Freshness lifetime and current age that \texttt{pingora-cache} evaluates from a
  retrieved object, at a given evaluation instant, shall match those it
  evaluates at that same instant from the object as it was handed to the
  backend --- including across a restart, and after a metadata update
  (\reqref{FRE-COMPLIANCE.REVALIDATION-01}), when the updated object is the
  one handed over.

  Two kinds of input feed that evaluation. The response header fields ---
  \texttt{Date}, \texttt{Expires}, \texttt{Age}, and the
  \texttt{Cache-Control} \texttt{max-age}/\texttt{s-maxage} directives --- are
  already required to round-trip byte for byte by
  \reqref{FRE-COMPLIANCE.HEADERS-01}, boundary and degenerate values included;
  interpreting them is up to \texttt{pingora-cache}. What this requirement
  adds is the other kind: the created, updated, and fresh-until instants
  \texttt{pingora-cache} keeps in the object's internal metadata, which the
  backend round-trips as part of the metadata it persists, never substituting
  timestamps of its own. Boundary values --- an \texttt{Age} of \texttt{0},
  delta-seconds at the overflow cap (\texttt{2147483648}), \texttt{max-age=0},
  an already-past or syntactically invalid \texttt{Expires} --- are where a
  mismatch in evaluation would surface, so this requirement's test exercises
  them.

  \textit{Basis: \rfcsec{9111}{4.2}[freshness], \rfcsec{9111}{4.2.3}[calculating age],
  \rfcsec{9111}{5.1}[Age], \rfcsec{9111}{5.2.2}[response directives],
  \rfcsec{9111}{5.3}[Expires].}

  \reqdiaglink{overview}{\textit{View the domain overview.}}

  \reqcover{verify}{TC}
  \reqcover{satisfy}{SDS}

  \deriveReqt{HLR-COMPLIANCE.RFC-01}
  \trace{FRE-COMPLIANCE.HEADERS-01}
  \trace{FRE-COMPLIANCE.REVALIDATION-01}
\end{requirement}

\begin{requirement}{FRE-COMPLIANCE.VALIDATORS-01}[draft][functionalRequirement]{Preserve validator fields}

  \localreqdiagram*{rfr05}{\reqid}{\reqdesc}

  The storage backend shall return the \texttt{ETag} and \texttt{Last-Modified}
  validator fields byte for byte in each form conditional revalidation depends
  on: a strong entity tag, a weak one (\texttt{W/}), an empty entity tag
  (\texttt{ETag: ""}), and both validators on one response.

  \reqref{FRE-COMPLIANCE.HEADERS-01} already requires every header field to
  round-trip byte for byte. This requirement names the validator forms whose
  silent normalisation would specifically break revalidation --- the fields
  \texttt{pingora-cache} uses to issue \texttt{If-None-Match} /
  \texttt{If-Modified-Since} and to match a \texttt{304 (Not Modified)} to the
  stored representation --- so that each is verified explicitly. Validators
  changed by a metadata update are governed by
  \reqref{FRE-COMPLIANCE.REVALIDATION-01}.

  \textit{Basis: \rfcsec{9110}{8.8.2}[Last-Modified], \rfcsec{9110}{8.8.3}[ETag];
  \rfcsec{9111}{4.3}[validation].}

  \reqdiaglink{overview}{\textit{View the domain overview.}}

  \reqcover{verify}{TC}
  \reqcover{satisfy}{SDS}

  \deriveReqt{HLR-COMPLIANCE.RFC-01}
  \trace{FRE-COMPLIANCE.HEADERS-01}
  \trace{FRE-COMPLIANCE.REVALIDATION-01}
\end{requirement}

\begin{requirement}{FRE-COMPLIANCE.REVALIDATION-01}[draft][functionalRequirement]{Freshen metadata}

  \localreqdiagram*{rfr06}{\reqid}{\reqdesc}

  WHEN \texttt{update\_meta} is called with newer metadata for the object a key
  holds, the storage backend shall replace that object's metadata and header
  fields with the newer set, keeping the object's body unchanged and
  retrievable under the same key, or, if the newer set exceeds the space the
  backend reserved for the object's metadata, stop holding the object, so
  that the next request for it fetches it again.

  \texttt{pingora-cache} freshens a stored response this way after a
  successful validation --- a \texttt{304 (Not Modified)} response --- through
  the backend's \texttt{update\_meta}
  operation.\footnote{\pgsrc{pingora-cache/src/storage.rs}{235}{\texttt{Storage::update\_meta}}.}
  A newer header set may be larger than the one it replaces, and no longer
  fit the space the on-disk layout reserved ahead of the body
  (\reqref{FDS-STORAGE.METAUPDATE-01}). For this version, the backend then
  stops holding the object rather than keep it
  (\reqref{ADR-STORAGE.METAMOVE-01}). \rfcsec{9111}{4.3.4} binds the
  responses a cache keeps, but nothing obliges a cache to keep one, and
  \rfcsec{9111}{4.4} counts removing a stored response as invalidating it.
  Either way, the object is never served truncated, nor with a mixture of
  the two metadata sets.
  The newer metadata belongs to the object that was validated: a key whose
  object a newer one has replaced since then keeps the newer object, whole,
  and the update is not applied (\reqref{FDS-STORAGE.METAUPDATE-01}).
  What a concurrent lookup may observe is stated in
  \reqref{FRE-COMPLIANCE.METAATOMIC-01}, what survives a restart or crash in
  \reqref{FRE-COMPLIANCE.METAPERSIST-01}, and what an update of a key with no
  stored object does in \reqref{FRE-COMPLIANCE.METAABSENT-01}. An update
  requested after the backend's process has received the Pingora upgrade signal
  is refused instead (\reqref{FRE-COMPAT.HOTUPGRADE-01}).

  \textit{Basis: \rfcsec{9111}{3.2}[updating stored header fields],
  \rfcsec{9111}{4.3.4}[freshening stored responses upon validation].}

  \reqdiaglink{overview}{\textit{View the domain overview.}}

  \reqcover{verify}{TC}
  \reqcover{satisfy}{FDS}

  \deriveReqt{HLR-COMPLIANCE.RFC-01}
  \trace{FDS-STORAGE.METAUPDATE-01}
\end{requirement}

\begin{requirement}{FRE-COMPLIANCE.METAATOMIC-01}[draft][functionalRequirement]{Atomic metadata update}

  \localreqdiagram*{rfrmetaatomic01}{\reqid}{\reqdesc}

  WHILE a metadata update of an object is in progress, the storage backend
  shall answer every lookup of that object with either its pre-update or its
  post-update metadata, whole, together with the object's unchanged body.

  A lookup therefore never observes a partially updated mixture of the two
  metadata sets, nor a body that no longer matches the metadata it is served
  with; a mixture could pair freshness inputs and validators that never
  belonged together. The update itself is stated in
  \reqref{FRE-COMPLIANCE.REVALIDATION-01}.

  \textit{Basis: \rfcsec{9111}{3.2}[updating stored header fields].}

  \reqdiaglink{overview}{\textit{View the domain overview.}}

  \reqcover{verify}{TC}
  \reqcover{satisfy}{FDS}

  \deriveReqt{HLR-COMPLIANCE.RFC-01}
  \trace{FRE-COMPLIANCE.REVALIDATION-01}
\end{requirement}

\begin{requirement}{FRE-COMPLIANCE.METAPERSIST-01}[draft][functionalRequirement]{Freshened metadata persists}

  \localreqdiagram*{rfrmetapersist01}{\reqid}{\reqdesc}

  WHEN a process of the backend looks up an object whose metadata update an
  earlier process reported as applied, the storage backend shall return the
  updated metadata --- or, only if a host crash or power loss intervened,
  either the updated or the superseded metadata, never a mixture of the two.

  Superseded metadata only makes the object due for revalidation sooner,
  whereas a mixture could pair freshness inputs and validators that never
  belonged together. The requirement holds until a later store, update, or
  purge of the key changes the object again.

  This requirement derives from two parents: \reqref{HLR-COMPLIANCE.RFC-01}
  for freshening semantics, which an update lost on an orderly restart would
  undo, and \reqref{HLR-PERSISTENCE.CRASHSAFE-01} for what a host crash may
  and may not do to an applied update.

  \textit{Basis: \rfcsec{9111}{4.3.4}[freshening stored responses upon
  validation].}

  \reqdiaglink{overview}{\textit{View the domain overview.}}

  \reqcover{verify}{TC}
  \reqcover{satisfy}{FDS}

  \deriveReqt{HLR-COMPLIANCE.RFC-01}
  \deriveReqt{HLR-PERSISTENCE.CRASHSAFE-01}
  \trace{FRE-COMPLIANCE.REVALIDATION-01}
\end{requirement}

\begin{requirement}{FRE-COMPLIANCE.METAABSENT-01}[draft][functionalRequirement]{Update of an absent key}

  \localreqdiagram*{rfrmetaabsent01}{\reqid}{\reqdesc}

  IF \texttt{update\_meta} is called for a key that holds no stored object,
  THEN the storage backend shall report an error distinguishable from an
  applied update, leaving the key without a stored object.

  Materialising a bodyless object from the metadata alone would let a later
  lookup serve a response with no representation. This deliberately differs
  from the stock backend, which panics in that
  case.\footnote{\pgsrc{pingora-cache/src/memory.rs}{448}{\texttt{MemCache::update\_meta}}
  calls \texttt{panic!("no meta found")} when the key holds no object.}

  \textit{Basis: \rfcsec{9111}{3.2}[updating stored header fields].}

  \reqdiaglink{overview}{\textit{View the domain overview.}}

  \reqcover{verify}{TC}
  \reqcover{satisfy}{FDS}

  \deriveReqt{HLR-COMPLIANCE.RFC-01}
  \trace{FRE-COMPLIANCE.REVALIDATION-01}
\end{requirement}

\begin{requirement}{FRE-COMPLIANCE.REPLACE-01}[draft][functionalRequirement]{Replace stored responses atomically}

  \localreqdiagram*{rfr10}{\reqid}{\reqdesc}

  The storage backend shall answer every lookup it serves from a published
  object with the object most recently published under the looked-up key,
  whole --- never a mixture of two objects, nor the metadata of one paired
  with the body of another.

  \texttt{pingora-cache} stores a new response under a key that already holds
  one when a full response to a validation request replaces the stored
  response it could not freshen, or when an expired object is simply
  re-fetched and re-admitted. The old object therefore stays the answer, whole,
  until the new store is published, and the new object is the answer, whole,
  from then on. Until then a lookup within the writing process may instead be
  served a partial hit of the in-flight new object
  (\reqref{FRE-COMPAT.STREAMING-01}); a lookup that cannot see the in-flight
  write --- such as one from a peer process --- gets the old object. When
  several stores of the same key race --- which \texttt{pingora-cache}
  permits (\reqref{FRE-COMPAT.CONCURRENT-01}) --- the last to publish is the
  single winner. What a reader already streaming the old object receives is
  stated in \reqref{FRE-COMPLIANCE.READERISOLATION-01}.

  \textit{Basis: \rfcsec{9111}{4.3.3}[handling a validation response --- a full
  response replaces the stored one]; \rfcsec{9111}{3}[storing responses].}

  \reqdiaglink{overview}{\textit{View the domain overview.}}

  \reqcover{verify}{TC}
  \reqcover{satisfy}{FDS}

  \deriveReqt{HLR-COMPLIANCE.RFC-01}
  \trace{FDS-STORAGE.PUBLISH-01}
  \trace{FRE-COMPAT.STREAMING-01}
\end{requirement}

\begin{requirement}{FRE-COMPLIANCE.READERISOLATION-01}[draft][functionalRequirement]{Readers survive replace/purge}

  \localreqdiagram*{rfrreaderisolation01}{\reqid}{\reqdesc}

  WHEN the object a reader is streaming is replaced or purged, that reader
  shall deliver only that object's own bytes, ending either at that object's
  normal end of body or with an error.

  A reader never crosses into the bytes of the replacing object, of another
  object, or of a torn record. Delivering only the object's own bytes implies
  that the storage the reader reads from is not reused for anything else while
  the reader still references it; the design decides how
  (\reqref{FDS-STORAGE.PUBLISH-01}, \reqref{FDS-STORAGE.PURGE-01}).

  \textit{Basis: \rfcsec{9111}{4.4}[invalidating stored responses];
  \rfcsec{9111}{3}[storing responses].}

  \reqdiaglink{overview}{\textit{View the domain overview.}}

  \reqcover{verify}{TC}
  \reqcover{satisfy}{FDS}

  \deriveReqt{HLR-COMPLIANCE.RFC-01}
  \trace{FRE-COMPLIANCE.REPLACE-01}
  \trace{FRE-COMPLIANCE.PURGE-01}
\end{requirement}

\begin{requirement}{FRE-COMPLIANCE.VARY-01}[draft][functionalRequirement]{Store and select Vary variants}

  \localreqdiagram*{rfr07}{\reqid}{\reqdesc}

  WHEN \texttt{pingora-cache} looks up a key, the storage backend shall return
  the object stored under exactly that key --- primary slot or variant slot
  --- or report a miss when none is stored there.

  The backend thus holds several response variants for the same primary cache
  key, each under the variance-augmented cache key \texttt{pingora-cache}
  presents, and never serves one slot's object in place of another's, so that
  \texttt{Vary}-based content negotiation is honoured. The variance is derived
  from the request fields named by the response's \texttt{Vary} header and
  folded into the cache key \emph{before} the backend is called; the backend
  never sees the request itself. Variants persist across a restart like any
  other stored object (\reqref{HLR-PERSISTENCE.DURABLE-01}); how mutating one
  variant affects the others is stated in
  \reqref{FRE-COMPLIANCE.VARYISOLATION-01}.

  \textit{Basis: \rfcsec{9111}{4.1}[calculating cache keys with Vary];
  \rfcsec{9110}{12.5.5}[Vary].}

  \reqdiaglink{overview}{\textit{View the domain overview.}}

  \reqcover{verify}{TC}
  \reqcover{satisfy}{SDS}

  \deriveReqt{HLR-COMPLIANCE.RFC-01}
\end{requirement}

\begin{requirement}{FRE-COMPLIANCE.VARYISOLATION-01}[draft][functionalRequirement]{Variant isolation}

  \localreqdiagram*{rfrvaryisolation01}{\reqid}{\reqdesc}

  WHEN one variant of a primary key is stored, replaced, freshened, or
  purged, the storage backend shall leave every other variant of that primary
  key unchanged.

  A purge addressed to the \emph{primary} slot is the deliberate exception: it
  makes every variant unreachable (\reqref{FRE-COMPLIANCE.PURGE-01}).

  \textit{Basis: \rfcsec{9111}{4.1}[calculating cache keys with Vary].}

  \reqdiaglink{overview}{\textit{View the domain overview.}}

  \reqcover{verify}{TC}
  \reqcover{satisfy}{SDS}

  \deriveReqt{HLR-COMPLIANCE.RFC-01}
  \trace{FRE-COMPLIANCE.VARY-01}
  \trace{FRE-COMPLIANCE.PURGE-01}
\end{requirement}

\begin{requirement}{FRE-COMPLIANCE.PURGE-01}[draft][functionalRequirement]{Honour purge and invalidation}

  \localreqdiagram*{rfr08}{\reqid}{\reqdesc}

  WHEN the storage backend applies an invalidating purge of a key, it shall
  serve neither the response stored under that key, nor any store of that
  key in flight when the purge arrives, nor --- for a primary key --- any of
  its \texttt{Vary} variants on any later lookup, until a store that begins
  after the purge publishes a new object under that key.

  The purge request path \texttt{pingora-proxy} drives addresses a purge to
  the \emph{primary} slot of an asset, and the variant lookups in
  \texttt{pingora-cache} consult the primary slot first; removing the primary slot therefore
  acts as the tombstone of the whole variant set, and the backend is not
  handed --- and cannot derive --- the keys of the individual variants. Such
  orphaned variant records stay in storage, counted against the cache's size
  budget like any other object, until eviction removes them
  (\reqref{FRE-COMPAT.EVICTION-01}). A same-key store in flight
  when the purge arrives may carry the very content the purge removes, if it
  began before the change that prompted the purge reached the origin, so the
  purge ends it as well: later lookups pass over it, and it publishes nothing
  when it completes. The request writing it still reads its own body back,
  and a reader already streaming it completes, as one streaming a purged
  record does (\reqref{FRE-COMPLIANCE.READERISOLATION-01}).
  The \texttt{pingora-cache} contract would let such a store
  land,\footnote{\pgsrc{pingora-cache/src/storage.rs}{217}{\texttt{Storage::expire}}'s
  doc calls a store that commits after an expiry ``the same race
  \texttt{Storage::purge} has'', to be ``resolved by whatever serializes
  writes for that key''.} but its in-memory backend removes a key's stores in
  flight together with its stored
  response,\footnote{\pgsrc{pingora-cache/src/memory.rs}{403}{\texttt{MemCache::purge}}
  removes the key from \texttt{self.temp}, its stores in flight, as well as
  from \texttt{self.cached}.} and an invalidation is the case where landing
  is wrong. A purge made to evict leaves stores in flight alone: it is not an
  invalidation. \texttt{pingora-cache} itself invalidates only on an
  explicit purge request from the proxy application --- it does not
  invalidate on unsafe request methods, leaving that RFC~9111 duty to the
  application.

  A purge of a key holding nothing is stated in
  \reqref{FRE-COMPLIANCE.PURGEABSENT-01}, when a purge may be acknowledged in
  \reqref{FRE-COMPLIANCE.PURGEDURABLE-01}, and what a reader streaming the
  purged object receives in \reqref{FRE-COMPLIANCE.READERISOLATION-01}.

  \textit{Basis: \rfcsec{9111}{4.4}[invalidating stored responses];
  \rfcsec{9110}{9.2.1}[safe methods].}

  \reqdiaglink{overview}{\textit{View the domain overview.}}

  \reqcover{verify}{TC}
  \reqcover{satisfy}{FDS}

  \deriveReqt{HLR-COMPLIANCE.RFC-01}
\end{requirement}

\begin{requirement}{FRE-COMPLIANCE.PURGEABSENT-01}[draft][functionalRequirement]{Purge or expiry of absent key}

  \localreqdiagram*{rfrpurgeabsent01}{\reqid}{\reqdesc}

  IF a purge or an expiry addresses a key under which no object is stored,
  THEN the storage backend shall report that nothing was found, with no error
  and no change to any other stored object.

  The case covers a key never stored and a key purged or expired a second
  time, so both operations are idempotent. \reqref{FRE-COMPAT.EXPIRE-01}
  states what an expiry does when the key \emph{does} hold an object.

  \textit{Basis: \rfcsec{9111}{4.4}[invalidating stored responses].}

  \reqdiaglink{overview}{\textit{View the domain overview.}}

  \reqcover{verify}{TC}
  \reqcover{satisfy}{FDS}

  \deriveReqt{HLR-COMPLIANCE.RFC-01}
  \trace{FRE-COMPLIANCE.PURGE-01}
  \trace{FRE-COMPAT.EXPIRE-01}
\end{requirement}

\begin{requirement}{FRE-COMPLIANCE.PURGEDURABLE-01}[draft][functionalRequirement]{Purge acknowledged once durable}

  \localreqdiagram*{rfrpurgedurable01}{\reqid}{\reqdesc}

  The storage backend shall report the outcome of an invalidating purge ---
  an object removed, or none found --- only once the key's absence will
  survive a crash or restart of the backend's process or host.

  An invalidation the backend will not make durable --- as after its process
  has received the Pingora upgrade signal (\reqref{FRE-COMPAT.HOTUPGRADE-01})
  --- therefore fails with an error rather than being reported as having
  removed something or nothing. An invalidation that found nothing is no
  exception: the key's record may have been removed just before by a purge
  made to evict, which nothing has made durable, and a crash undoing that
  removal would bring the object back after the invalidation reported it
  absent.

  A purge made to evict (\texttt{PurgeType::Eviction}) is outside this
  requirement. No one asked for the object's removal: the cache chose to
  drop it for space, and an evicted object back after a crash is still a
  valid response, which the cache may serve or evict again. It is not a
  purge in \reqref{HLR-PERSISTENCE.CRASHSAFE-01}'s sense.

  This requirement derives from two parents: \reqref{HLR-COMPLIANCE.RFC-01}
  for invalidation semantics, which a purged object reappearing would undo,
  and \reqref{HLR-PERSISTENCE.CRASHSAFE-01} for the durability of an
  acknowledged purge across a host crash.

  \textit{Basis: \rfcsec{9111}{4.4}[invalidating stored responses].}

  \reqdiaglink{overview}{\textit{View the domain overview.}}

  \reqcover{verify}{TC}
  \reqcover{satisfy}{FDS}

  \deriveReqt{HLR-COMPLIANCE.RFC-01}
  \deriveReqt{HLR-PERSISTENCE.CRASHSAFE-01}
  \trace{FRE-COMPLIANCE.PURGE-01}
\end{requirement}

\begin{requirement}{FRE-COMPLIANCE.INCOMPLETE-01}[draft][functionalRequirement]{Never serve incomplete responses}

  \localreqdiagram*{rfr09}{\reqid}{\reqdesc}

  IF a store does not complete --- it aborts, or its process or host crashes
  first --- or a host crash occurs before a completed store's content has
  reached stable storage, THEN the storage backend shall serve that store's
  object on no later lookup, in the writing process or any later one.

  A lookup that would have returned the object reports a miss instead, or ---
  when the key already held an object --- returns that earlier object
  (\reqref{FRE-COMPLIANCE.REPLACE-01}); truncated, torn, or otherwise
  corrupted content is never served as complete, and nothing such a store
  leaves on disk is promoted by a later process. The \texttt{Storage}
  interface gives \texttt{pingora-cache} no way to consume an ``incomplete''
  record, so a miss is the only compliant outcome. The backend does not
  combine partial content from \texttt{206 (Partial Content)} or range
  responses, so it stores only complete representations and leaves any range
  handling to \texttt{pingora-cache}. What a reader already following the
  store receives when it aborts is stated in
  \reqref{FRE-COMPLIANCE.PARTIALABORT-01}, and the release of the storage
  such a store consumed in \reqref{FRE-COMPAT.ABORTSPACE-01} and
  \reqref{FRE-COMPAT.STARTRECLAIM-01}.

  This requirement derives from two parents: \reqref{HLR-COMPLIANCE.RFC-01}
  for the RFC~9111 rule on storing incomplete responses, and
  \reqref{HLR-PERSISTENCE.CRASHSAFE-01} for extending it to content a host
  crash leaves incomplete.

  \textit{Basis: \rfcsec{9111}{3.3}[storing incomplete responses],
  \rfcsec{9111}{3.4}[combining partial content]; \rfcsec{9110}{14.2}[Range],
  \rfcsec{9110}{6.1}[framing and completeness].}

  \reqdiaglink{overview}{\textit{View the domain overview.}}

  \reqcover{verify}{TC}
  \reqcover{satisfy}{FDS}

  \deriveReqt{HLR-COMPLIANCE.RFC-01}
  \deriveReqt{HLR-PERSISTENCE.CRASHSAFE-01}
\end{requirement}

\begin{requirement}{FRE-COMPLIANCE.PARTIALABORT-01}[draft][functionalRequirement]{Aborted partial hits end in error}

  \localreqdiagram*{rfrpartialabort01}{\reqid}{\reqdesc}

  IF the in-flight store a partial hit's reader is attached to aborts, THEN
  that reader shall end its stream with an error, never with a normal end of
  body.

  A truncated representation is thus never delivered to a follower as if it
  were complete. The abort may be an upstream failure, a cancellation, or a
  write failure (\reqref{FRE-COMPAT.WRITEFAIL-01}); the partial hit itself is
  stated in \reqref{FRE-COMPAT.STREAMING-01}, and its normal end on
  completion in \reqref{FRE-COMPAT.STREAMEND-01}.

  \textit{Basis: \rfcsec{9111}{3.3}[storing incomplete responses];
  \rfcsec{9110}{6.1}[framing and completeness].}

  \reqdiaglink{overview}{\textit{View the domain overview.}}

  \reqcover{verify}{TC}
  \reqcover{satisfy}{FDS}

  \deriveReqt{HLR-COMPLIANCE.RFC-01}
  \trace{FRE-COMPLIANCE.INCOMPLETE-01}
  \trace{FRE-COMPAT.STREAMING-01}
\end{requirement}

